\section{Introduction}

In this paper, we present the classification
of the positive harmonic functions
on the discrete Heisenberg group $G=H_3(\m Z)$.

\subsection{Positive harmonic functions}

Let $\mu=\sum_{s\in S} \mu_s \de_s$ be a positive measure on $G$
with finite support $S\subset G$.
We recall that a function $h$ on $G$ is said to be $\mu$-harmonic
if it satisfies the equality
$h= P_\mu h$, where $P_\mu h(g):= \sum_{s\in S} \,\mu_s\, h(sg)$ for all $g$ in $G$.
We~want to describe the cone $\mc H_\mu^+$ of positive $\mu$-harmonic functions $h$ on $G$.
By the Choquet theorem, it is enough to describe its extremal rays.

The main aim of this paper is to prove
that the extremal positive $\mu$-harmonic functions
on $G$ are proportional
either to a character of $G$ or to a translate of a function which is induced
from a character of an abelian subgroup (Theorem \ref{thharmu1}).

The special case
where $\mu$ is the {\it southwest measure} was handled in
the introductory paper
\cite{BenoistHeisenberg1}.
This case was striking because the classical partition
function $h(x,y,z):=p_y(z)$ with
\[
p_y(z):=
\text{number of partition of $z$
by $y$ non-negative integers}
\]
occurs as one of these extremal
positive harmonic functions.
This partition function $p_y(z)$ is the simplest instance
of a ``harmonic function induced
from the character of an abelian subgroup''
that we will introduce in this paper.

\subsection{Construction of harmonic functions}
\label{secincoha}
The simplest examples of $\mu$-harmonic functions are $\mu$-harmonic characters.
Those are the
characters
$\chi: G\ra \m R_{>0}$ such that $\sum_{s\in S}\mu_s\,\chi(s)=1$.
Such a function
$h=\chi$ is an
extremal positive $\mu$-harmonic function on $G$ which is invariant by the center $Z$ of $G$,
see Lemma \ref{lemcarext}.

We now introduce another construction of extremal positive $\mu$-harmonic functions
by inducing harmonic characters.
Let $S_0\subset S$ be a maximal abelian subset and $G_0$
be the subgroup of $G$ generated by $S_0$. Denote by
$\mu_{0}:=\sum_{s\in G_{_0}}\mu_s\, \de_s$
the measure restriction of $\mu$ to $G_0$.
Let $\chi_0$ be a $\mu_{0}$-harmonic character of $G_{0}$.
We extend $\chi_0$ as a function on $G$,
still denoted $\chi_0$,
which is $0$ outside $G_{0}$. This function $\chi_0$ is
$\mu$\nobreakdash-subharmonic, so that the sequence $P_\mu^n\chi_0$
is increasing. We set
\[
h_{G_{_0},\chi_{_0}}=\lim_{n\ra\infty}P_\mu^n\chi_0.
\]
We will tell exactly for which pairs $(G_0,\chi_0)$
the function $h_{G_{_0},\chi_{_0}}$ is finite,
in Lemma~\ref{lemindcha} and in
Propositions \ref{proindra1}, \ref{proindra2} and \ref{proindra3}.
When it is finite, the function $h_{G_{_0},\chi_{_0}}$
is an
extremal positive $\mu$-harmonic function on $G$,
see Lemma \ref{lemconind}.
We will call $h_{G_{_0},\chi_{_0}}$
the \emph{harmonic function on $G$
induced from the $\mu_0$-harmonic character $\chi_0$ of $G_0$}.

For $g$ in $G$, we denote by $\rho_{g}:g'\mto g'g$
the right translation by $g$ on $G$.
Whenever a function $h$ is $\mu$-harmonic,
the function $h_{g}:=h\circ\rho_{g}$
is also $\mu$-harmonic.

\subsection{Main results}
\label{secfinsup}

Our main theorem tells us that conversely
these three constructions are the only possible ones.

\begin{Thm}
\label{thharmu1}
Let $G=H_3(\m Z)$ be the discrete Heisenberg group
and $\mu$ be a positive measure on $G$
whose support $S$ is finite and generates the group $G$.
Then every extre\-mal positive $\mu$-harmonic function $h$
on $G$ is proportional either to a character $\chi$ of $G$ or
to a translate $h_{G_{_0},\chi_{_0}}\circ \rho_{g}$ of a function
induced from a harmonic character of an abelian subgroup.
\end{Thm}

\begin{figure}[htb]
\centerline{\includegraphics[scale=.6]{case1.pdf}}
\caption{\label{figcase1}In Case \eqref{thconclu1} and in Case \eqref{thconclu2b} of Theorem \ref{thconclu}, no harmonic function is induced from a character
of an abelian subgroup~$G_0$.}
\end{figure}

\begin{figure}[htb]
\centerline{\includegraphics[scale=.65]{case2.pdf}}
\caption{\label{figcase2}In Case \eqref{thconclu2a}, exactly two harmonic functions are induced from a character of $G_0\! =\! G_{\mu_{_0}}$ and no other.
In Case \eqref{thconclu3a}, only one harmonic function is induced
from a character of $G_0\!=\! G_{\mu_{_0}}$ and one or infinitely many are induced
from a character of $G_1\!=\! G_{\mu_{_1}}$.}
\end{figure}

\begin{figure}[htb]
\centerline{\includegraphics[scale=.65]{case3.pdf}}
\caption{\label{figcase3}In case \eqref{thconclu3b}, infinitely many harmonic functions are induced from a character of $G_0=G_{\mu_{_0}}$ and one or infinitely many are induced from a character of $G_1=G_{\mu_{_1}}$.}
\end{figure}

\skpt
\begin{Rem}
\begin{itemize}
\item
Of course the case where $\mu(G)=1$
is the major case.
However, even when dealing with a probability measure $\mu$,
the induction process forces us to work with positive measures
$\mu_0$ which are not probability measures.

\item
Theorem \ref{thharmu1} can not be extended to all
nilpotent groups $G$. Indeed, the conclusion of Theorem
\ref{thharmu1} is not always valid for
a probability measure $\mu$ on the
nilpotent group $G$ of rank 4 with cyclic center.
See Section \ref{secrankfour}.
\end{itemize}
\end{Rem}

Theorem \ref{thharmu1} has been announced in
\cite{BenoistHeisenberg1}.
It will be proved in Section \ref{secinvhar}. Indeed it is a
direct consequence of Propositions \ref{provmulin1} and \ref{provmunol}.
We will give a more precise description of the
extremal positive $\mu$-harmonic functions $h$
in Theorem \ref{thconclu}.
In particular, we~will say exactly when and how many of these new examples occur. This is illustrated in the schematic Figures
\ref{figcase1}, \ref{figcase2} and \ref{figcase3}. In these figures, we have drawn
various cases of semigroup $G_\mu^+$ generated by $S$
that are described in Theorem \ref{thconclu}.
Note that the support of a positive $\mu$-harmonic function $h$ is invariant by the opposite semigroup, \ie by the semigroup generated by $S^{-1}$.
In particular when $G_\mu^+=G$, a~positive harmonic function $h$ is either
identically zero or vanishes nowhere.
Here are two corollaries of Theorem \ref{thconclu} that we will prove in Section \ref{secconclu}.
The first corollary tells us that these new examples always
vanish somewhere.\enlargethispage{.5\baselineskip}%

\begin{Cor}
\label{corharmu1}
Same notation.
Let $h$ be an extremal positive $\mu$-harmonic function
on $G$ which does not vanish. Then $h$ is a character of $G$.
\end{Cor}

The second corollary tells us exactly when no new example occurs.
We denote by~$G_\mu^+$ the semigroup generated by $S$.

\begin{Cor}
\label{corharmu2}
Same notation with $\mu(G)=1$.
The following are equivalent:
\begin{enumeratei}
\item\label{corharmu2i}
Every extremal positive $\mu$-harmonic function $h$ on $G$ is
a character of $G$.
\item\label{corharmu2ii}
$G_\mu^+$ contains two non-central elements whose product is in
$Z\smallsetminus \{0\}$.
\end{enumeratei}
\end{Cor}

\subsection{Previous results}
\label{secpreres}
\bq
The study of harmonic functions on groups has a very long history. I will just point out the part of it which is relevant
for our purposes.
\eq

As a general motivation, let us recall that
the bounded $\mu$-harmonic functions on a group $G$
are described thanks
to bounded functions on the \emph{Poisson boundary} of $(G,\mu)$.
They are used to study
random walks on $G$-spaces.
The extremal positive $\mu$\nobreakdash-harmonic functions on $G$
are related to the Martin boundary of $(G,\mu)$.
They are used to study
more precisely the behavior of these random walks,
see \cite{Ancona90}, \cite{Gouezel15}, \cite{Sawyer97} or~\cite{Woess00}.

\subsubsection{Abelian groups}
This part of the history begins with the Choquet--Deny theorem
in \cite{ChoquetDeny60}:

\emph{Let $G$ be a finitely generated abelian group and
$\mu$ be a
positive finite measure on~$G$
whose support generates $G$ as a group.
Then every extremal positive $\mu$-harmonic function
$h$ on $G$ is proportional to a character.}

Indeed the proof of this theorem is very short:
one notices that the harmonicity equation \eqref{eqnfghar}
is a decomposition of $h$ as a sum of positive harmonic functions
and hence all the terms in this sum are proportional to $h$.

\subsubsection{Bounded harmonic functions}
The Choquet--Deny theorem has been extended to nilpotent
groups when $\mu$ has mass $1$
and $h$ is bounded. This is due to
Dynkin and Maljutov in \cite{DynkinMaljutov61}:

\emph{Let $G$ be a finitely generated nilpotent group
and $\mu$ be a probability measure on $G$
whose support generates $G$ as a group.
Then every bounded $\mu$-harmonic
function on~$G$ is constant.}

\subsubsection{When $S$ generates $G$ as a semigroup}
The Choquet--Deny theorem has also been extended to nilpotent
groups for $h$ unbounded under an extra assumption.
This is due to Margulis in \cite{Margulis66}:

\emph{Let $G$ be a finitely generated nilpotent group
and $\mu$ be a positive measure on $G$
whose support generates $G$ \emph{as a semigroup}.
Then every extremal positive $\mu$-harmonic
function on $G$ is proportional to a character.}

\subsubsection{The Heisenberg group}
The main significance of our Theorem \ref{thharmu1} is
that even though the Choquet--Deny theorem can not be
extended to finitely generated nilpotent groups
without this extra assumption, for the Heisenberg group
one can describe all the positive harmonic functions.
Note that, because of Margulis theorem,
most of our paper will deal
with a positive measure whose support
generates $G$ as a group but does not
necessarily generate $G$ as a semigroup.

Many recent works focus
on the random walks on the discrete Heisenberg group
$G$ as in
\cite{Breuillard05},
\cite{DiaconisHough15} and
\cite{GollSchmidtVerbitskiy16},
or on nilpotent groups as in \cite{Breuillard10}
and \cite{Guivarch71},
or on the geometry of words in $G$ as in
\cite{LindSchmidt15} and \cite{VershikMalyutin18}.
We mention these related results
even though we will not use them.

\subsection{Strategy of proof}
\label{secstrpro}
\bq
We now explain the strategy of proof of Theorem \ref{thharmu1}
and the organization of the paper.
\eq

In Section \ref{secnotpre}, we recall well-known facts on positive harmonic functions
and notations for the discrete Heisenberg group $G$
and its positive measures $\mu$ with a finite support $S$.

In Section \ref{secproof}, we begin the proof of Theorem \ref{thharmu1}.
When $h$ is an extremal $\mu$-harmonic function on $G$, we merely focus on the
equality
$h(g)=P^n_\mu h(g)$, where the right-hand side is written as a weighted sum
of values $h(\dot{w}g)$ for words $w$ of length $n$ in $S$,
as in Equation \eqref{eqnpmun}.
In Lemmas \ref{lemconind} and \ref{lemrecind},
we check that when the contribution in this sum
of the words $w$ whose letters are in a proper subgroup of $G$,
is not negligible,
then~$h$ is an ``induced harmonic function''.
In Lemma \ref{lemnegtra}, we prove a useful generalization:
we allow $w$ to be a concatenation of $k$ subwords
whose letters are in a proper subgroup with $k \geq 1$ fixed.
The proofs are very general
and do not assume $G$ to be nilpotent.

In Section \ref{secinvhar}, we assume that ``$h$ is not induced'',
and we want to prove that $h$ is invariant by the center $Z$ of $G$.
The main idea is to construct a symmetric relation~$\mc R_n$
among the words in $S^n$ such that
two related words $w$ and $w'$ have same weight
and their image $\dot{w}$ and $\dot{w}'$ in $G$
differ by a non-trivial element $z$ of $Z$.
A key point is to be able to compare the number of words related to $w$
and the number of words related to $w'$, see Lemma \ref{lemkwkw}. This allows
us to prove that $h$ is proportional to one of its translate $h_z$,
see Proposition \ref{proseminv}.
The last step is to prove that $h$ is indeed equal to its translate $h_z$.
This is done in Propositions \ref{provmulin1} and \ref{provmunol}.
The key point there, Lemma~\ref{lemsumbww} is based
on a counting argument that again involves the partition function.
This finishes the proof of Theorem \ref{thharmu1}.

In Section \ref{secexiind}, we give a complete classification of the
extremal $\mu$-harmonic functions that are ``induced
from a character'',
see Theorem \ref{thconclu}.
Their existence is an important new feature of this article.
The proof of this classification in Propositions \ref{proindra1}, \ref{proindra2} and~\ref{proindra3}
uses a transience property for random walks on $\m Z$
similar to the large deviation inequality, see Lemma \ref{lemranwal}.

In the last Section \ref{secrankfour}, we explain how to construct,
for a rank $4$ nilpotent group,
new extremal positive $\mu$-harmonic
functions that are not induced.

\section{Notation and preliminary results}
\label{secnotpre}

We introduce in this section notations
that will be used all over
this article.

\subsection{The cone of $\mu$-harmonic functions}
\label{secconhar}

\bq
We first recall classical facts on positive $\mu$-harmonic functions.
\eq

Let $G$ be a finitely generated group and $\mu$ be a positive
measure with finite support $S\subset G$.
We denote by $G_\mu^+$ the subsemigroup of $G$ generated by $S$
and by $G_\mu$ the subgroup of $G$ generated by $S$.

A positive function $h:G\ra [0,\infty[$ is said to be
\emph{$\mu$-harmonic} if it satisfies the equality
\begin{equation}
\label{eqnfghar}
h= P_{\mu}h,
\quad\text{where}\quad
P_{\mu}h: g\mto \sum_{s\in S}\mu_s h(sg).
\end{equation}
A non-zero positive $\mu$-harmonic function is said to be \emph{extremal
or $\mu$-extremal} if every smaller positive $\mu$-harmonic function
$h'\leq h$ is a multiple of $h$.

A function $h$ is said to be \emph{$\mu$-superharmonic},
respectively \emph{$\mu$-subharmonic},
if it satisfies the inequality $h\geq P_{\mu}h$, respectively $h\leq P_{\mu}h$.

We will often write the $n^\mathrm{th}$ power of the operator $P_\mu$ under the form
\begin{equation}
\label{eqnpmun}
P_{\mu}^nh(g)= \sum_{w\in S^n}\mu_w h(\dot{w}g),
\end{equation}
where, for a word $w=s_1\ldots s_n\in S^n$ of length $\ell_w=n$,
the constant $\mu_w>0$ is the product
$\mu_w:=\mu_{s_{_1}}\cdots \mu_{s_{_n}}>0$
and where the element $\dot{w}\in G$ is the product $\dot{w}:=s_1\cdots s_n$ in $G$.

Let $\mc H_\mu^+$ be the convex cone of positive
$\mu$-harmonic functions $h$ on $G$
and $\mc E$ be a Borel set of extremal $\mu$-harmonic functions containing exactly one
function in each extremal ray of $\mc H_\mu^+$.
We endow $\mc H_\mu^+$ with the topology of the pointwise convergence.
When $G_\mu^+=G$ the cone $\mc H_\mu^+$ \emph{has a compact basis},
this means that there exists a compact subset of $\mc H_\mu^+$
that meets all rays of $\mc H_\mu^+$.
In general, the cone $\mc H_\mu^+$ might not have a compact basis but it is \emph{well-capped},
this means that it is a union of closed convex subcones $\mc H_{\mu,i}^+$
with compact basis such that $\mc H_\mu^+\smallsetminus \mc H_{\mu,i}^+$
is also convex.
This cone~$\mc H_\mu^+$ is also \emph{reticulated}, this means that
every two positive $\mu$-harmonic functions~$h_1$ and~$h_2$
admit a maximal $\mu$-harmonic lower bound $h_m$ and also
a minimal $\mu$-harmonic upper bound $h_M$.
Indeed one has
\begin{align*}
h_m
&=
\lim_{n\ra\infty}P^n_\mu (\min(h_1,h_2))\geq 0
\quad\text{and}\quad \\
h_M
&=
\lim_{n\ra\infty}P^n_\mu (\max(h_1,h_2))
\leq h_1+ h_2<\infty.
\end{align*}

By the Choquet theorem, it is enough to describe the extremal rays of this cone~$\mc H_\mu^+$.
Indeed, since $\mc H_\mu^+$ is well-capped, this theorem tells us
that \emph{every positive $\mu$-harmonic function $h$
can be written as an integral of non-proportional
extremal $\mu$-harmonic functions:
$h=\int_{\mc E}f\rd \al(f)$,
for a positive measure $\al$
on the set $\mc E$}.

Since $\mc H_\mu^+$ is reticulated,
this theorem also tells us that \emph{such a measure $\al$ is unique.}

In this paper a \emph{character}
will always mean a multiplicative morphism $\chi:G\mto \m R_{>0}$.
A character $\chi$ is $\mu$-harmonic if and only if it satisfies
the equation $\sum_{s\in S} \mu_s\,\chi(s)=1$.

\subsection{Harmonic characters}
\bq
We discuss here harmonic characters on nilpotent groups.
\eq

Let $G$ be a nilpotent finitely generated group
and $\mu$ be a positive finite measure on $G$
with finite support generating $G$.

\begin{Lem}
\label{lemcarext}
Every $\mu$-harmonic character of $G$ is an extremal
positive $\mu$-harmonic function.
\end{Lem}

\begin{proof}[Proof of Lemma \ref{lemcarext}]
Let $\chi$ be a $\mu$-harmonic character
such that $\chi=h'+h''$
with both $h'$ and $h''$
positive and $\mu$-harmonic.
We want to prove that the function $\widetilde{h}':=\chi^{-1}h'$
is constant.
We notice that the measure $\widetilde\mu:=\chi\mu$ on $G$
is a probability measure and the function $\widetilde h'$
is a bounded $\widetilde\mu$-harmonic function.
Therefore by Dynkin--Maljutov theorem,
see Section \ref{secpreres}, the function $\widetilde h'$ is constant.
\end{proof}

\subsection{The Heisenberg group}
\label{secheigro}

\bq
We gather here notation that we will use in this article
for the discrete Heisenberg group.
\eq

Recall that the discrete Heisenberg group $G:=H_3(\m Z)$ is the set $\m Z^3$
of triples seen as matrices
$(x,y,z):=\mbox{\scriptsize
$\left(\!
\begin{array}{ccc} 1 &x&z \\
0 &1&y\\
0&0&1
\end{array}\!
\right)$}.
$
It is endowed with the product
\begin{equation}
\label{eqnprohei}
(x,y,z)\, (x',y',z')= (x+x',y+y',z+z'+xy').
\end{equation}
We will denote by $0:=(0,0,0)$ the identity element of $G$,
and by $z_0$ the generator $z_0:=(0,0,1)$ of the center $Z$ of $G$.

For two elements $g=(x,y,z)$, $g'=(x',y',z')$ of $G$, we will denote
by $c_{g,g'}$ the integer $c_{g,g'}:=xy'\!-\!yx'$ so that
\begin{equation}
\label{eqncggzoc}
gg'g^{-1}{g'}^{-1}={z_0}^{c_{g,g'}}.
\end{equation}

Let $\ol{G}:=G/Z\simeq \m Z^2$ be the abelianization of $G$
that we embed in the real vector space
$V:=\ol{G}\otimes_{\m Z}\m R\simeq \m R^2$.

Let $\mu$ be a positive measure on $G$ with finite support $S$.
We denote by $\ol{\mu}$ the image of $\mu$ in $\ol{G}$ and
by $\ol{S}$ its support.

We denote by $V_\mu$ the vector subspace of $V$ generated by $\ol{S}$
and by $V^+_\mu$ the smallest convex cone of $V$ containing $\ol{S}$.
Note that, when $G_\mu=G$, one always has $V_\mu=V$,
and, when $G^+_\mu=G$, one always has $V^+_\mu=V$.

The description of $\mc H_\mu^+$, when $G_\mu=G$ will heavily depend on the shape of $V^+_\mu$.
We will often distinguish the three cases:
\begin{equation}
\label{eqnvmuphc}
\mbox{ $V^+_\mu$ = the plane, a half-plane, or a properly convex cone.}
\end{equation}

\section{Induced harmonic functions}
\label{secproof}

In this section we present general facts on $\mu$-harmonic functions
on a finitely generated group $G$. These facts will be particularly useful
when $G$ is the Heisenberg group.

\subsection{Construction of induced harmonic functions}
\label{secconind}

\bq
The following lemma gives us a method to construct
$\mu$-harmonic functions starting from a harmonic function
for a smaller measure $\mu_0$.
This lemma will be mainly useful when $\mu_0$ is the restriction of $\mu$
to a proper subgroup $G_0$.
\eq

Let $G$ be a finitely generated group and
$\mu$ and $\mu_0$ be positive measures on $G$ with finite support
such that $\mu_0 <\mu$, \ie such that $\mu_1:=\mu-\mu_0$ is also a positive measure.

\begin{Lem}
\label{lemconind}
Let $h_0$ be a positive $\mu_0$-harmonic function on $G$ such that
the function $h:=\sup_{n\geq 1}P^n_{\mu}h_0$ is finite.
\begin{enumeratei}
\item\label{lemconindi}
Then
one has
$h=\lim_{n\ra\infty}P^n_{\mu}h_0$ and
$h$ is a positive $\mu$-harmonic function.
\item\label{lemconindii}
One can recover $h_0$ from $h$ as
$h_0=\lim_{n\ra\infty}P^n_{\mu_{_0}}h$.
\item\label{lemconindiii}
Moreover when $h_0$ is $\mu_0$-extremal then $h$ is $\mu$-extremal too.
\end{enumeratei}
\end{Lem}

When it is finite, the function $h$ will be called \emph{induced from the harmonic function~$h_0$}.

\skpt
\begin{proof}[Proof of Lemma \ref{lemconind}]

\eqref{lemconindi}
We first notice that, since $h_0=P_{\mu_{_0}}h_0\leq P_\mu h_0$,
the sequence $P^n_\mu h_0$ is increasing.
Hence, when this sequence is bounded it converges to a $\mu$-harmonic function.

\eqref{lemconindii}
Since $h\!=\!P_{\mu}h\!\geq\!P_{\mu_{_0}} h$,
the sequence $P^n_{\mu_{_0}} h$ is decreasing.
Since \hbox{$P_{\mu_{_0}}^n h\!\geq\!P_{\mu_{_0}}^n h_0\!=\!h_0$}, this sequence $P^n_{\mu_{_0}} h$
converges to a $\mu_0$-harmonic function
$h'_0:=\lim_{n\ra\infty}P^n_{\mu_{_0}}h$ such that
$h'_0\geq h_0$.

We want to prove that the function $h''_0:= h'_0-h_0$ is zero.
Since $h_0\leq h'_0\leq h$, one has
$P^n_\mu h_0\leq P^n_\mu h'_0\leq h$.
Therefore one also has
$\lim_{n\ra\infty}P^n_{\mu}h'_0=h$
and hence
$\lim_{n\ra\infty}P^n_{\mu}h''_0=0$.
Since $h''_0$ is $\mu_0$-harmonic, this last sequence is increasing
and hence one has $h''_0=0$.

\eqref{lemconindiii}
Assume now that $h_0$ is $\mu_0$-extremal and assume that
$h$ is the sum of two positive $\mu$-harmonic functions
$h=h'+h''$. We want to prove that $h$ and $h'$ are proportional.
The functions
$h'_0=\lim_{n\ra\infty}P^n_{\mu_{_0}}h'$
and
$h''_0=\lim_{n\ra\infty}P^n_{\mu_{_0}}h''$
are $\mu_0$-harmonic and, by~\eqref{lemconindii}, they give a decomposition
$h_0=h'_0+h''_0$.

Therefore,
one has $h'_0=\la' h_0$ and $h''_0=\la'' h_0$
for positive constants $\la'$ and $\la''$ with $\la'+\la''=1$.
One has the inequalities
\[
h'\geq\lim_{n\ra\infty}P^n_{\mu}h'_0 =\la' h\quad\text{and}
\quad h''\geq\lim_{n\ra\infty}P^n_{\mu}h''_0 =\la'' h.
\]
Since $h=h'+h''$, these inequalities are equalities: one has
$h'=\la' h$ and $h''=\la'' h$.
This proves that the function $h$ is $\mu$-extremal.
\end{proof}

\subsection{Recognizing induced harmonic functions}
\label{secrecind}

\bq
The following lemma is a converse of Lemma \ref{lemconind}.
It tells us how to recognize
a $\mu$-harmonic function that is induced from a $\mu_0$-harmonic function.
\eq

Let $G$ be a finitely generated group and
$\mu_0<\mu$ be positive measures on $G$ with finite support.

\begin{Lem}
\label{lemrecind}
Let $h$ be a positive $\mu$-harmonic function on $G$ such that
the function $h_0:=\inf_{n\geq 1}P^n_{\mu_{_0}}h$ is non-zero.
\begin{enumeratei}
\item\label{lemrecindi}
Then
one has
$h_0=\lim_{n\ra\infty}P^n_{\mu_{_0}}h$ and
$h_0$ is a positive $\mu_0$-harmonic function.
\item\label{lemrecindii}
One has the inequality
$h\geq\lim_{n\ra\infty}P^n_{\mu}h_0$.
\item\label{lemrecindiii}
Moreover when $h$ is $\mu$-extremal, one has the equality
$h=\lim_{n\ra\infty}P^n_{\mu}h_0$ and~$h_0$ is $\mu_0$-extremal too.
\end{enumeratei}
\end{Lem}

In particular, when $h$ is $\mu$-extremal, $h_0$ is supported by a translate
$G_{\mu_{_0}}g$ of the subgroup $G_{\mu_{_0}}$.

\begin{proof}[Proof of Lemma \ref{lemrecind}]
The argument is very similar to that of Lemma \ref{lemconind}.

\eqref{lemrecindi}
Since the function $h$ is positive and $\mu$-harmonic, the sequence
$P^n_{\mu_{_0}}h$ is positive and decreasing. Hence it has a limit $h_0$
which is $\mu_0$-harmonic.

\eqref{lemrecindii}
By assumption, this limit $h_0$ is non-zero. By construction,
one has the inequality
$h\geq h_0$. Since $h$ is $\mu$-harmonic, the sequence
$ P^n_{\mu}h_0$ is bounded by $h$ and, by Lemma~\ref{lemconind},
the limit $h':=\lim_{n\ra\infty}P^n_\mu h_0$
exists, is $\mu$-harmonic and is bounded by $h$.

\eqref{lemrecindiii}
Assume now that $h$ is $\mu$-extremal.
Then one has $h'=\la' h$ for some constant $\la'\geq 0$.
Again by Lemma \ref{lemconind}, one also has
\begin{equation}
\label{eqnhlipnh}
h_0=\lim_{n\ra\infty}P^n_{\mu_{_0}}h'
=\la'\lim_{n\ra\infty}P^n_{\mu_{_0}}h=\la' h_0.
\end{equation}
Therefore one has $\la'=1$.

It remains to check that $h_0$ is $\mu_0$-extremal.
Assume that $h_0=h'_0+h''_0$ with both~$h'_0$ and~$h''_0$
positive $\mu_0$-harmonic.
The limit
$h'':=\lim_{n\ra\infty}P^n_{\mu}h''_0$
is a $\mu$-harmonic function bounded by $h$.
Hence one has $h''=\la'' h$
and by the same computation as \eqref{eqnhlipnh}, one gets
$h''_0=\la'' h_0$.
This proves that $h_0$ is extremal.
\end{proof}

The following definition relies on the previous lemmas:

\begin{Def}

\label{defnonind}
A $\mu$-harmonic function $h$ on $G$
is said to be \emph{induced from a subgroup $G_0$}
if
\begin{equation}
\label{eqnlimpmunh}
\lim_{n\ra\infty}P_{\mu_{_0}}^nh\neq 0,
\end{equation}
where $\mu_0$ is the restriction of $\mu$ to $G_0$.

By Lemma \ref{lemrecind}, when $h$ is $\mu$-extremal this limit
\eqref{eqnlimpmunh} is equal to $h_0\circ g$, where
$g$ is in $G$ and
$h_0$ is an extremal $\mu_0$-harmonic function supported on $G_0$.
Therefore one has $h=h_{G_{_0},h_{_0}}\circ \rho_g$,
where $h_{G_{_0},h_{_0}}:=
\lim_{n\ra\infty}P_{\mu}^nh_0$.
In this case \emph{the function $h$ is a translate of
the harmonic function induced from $h_0$}.
Equivalently, \emph{the function $h$ is induced from $h_0\circ \rho_g$}.
\end{Def}

\begin{Def}

\label{defnonindsuite}
A $\mu$-harmonic function is said to be {induced},
if there exists a subgroup $G_0$ of infinite index in $G$
such that $h$ is induced from $G_0$.
It is said to be \emph{non-induced} otherwise.
\end{Def}

\begin{Rem}

The reason why we require in this definition $G_0$ to have infinite index
will be explained in Lemma \ref{lemnonind}.
A posteriori, for an extremal positive $\mu$-harmonic function
$h$ on the Heisenberg group $G$ with $G_\mu=G$,
this requirement is not so useful.
Indeed, by Corollary \ref{corconind}, the characters of $G$ are not induced
from proper finite index subgroups.
Moreover, by Definition \ref{defnonind}, if $h$ is induced from an infinite index subgroup~$G_0$, it is also induced from all the finite index subgroup of $G$ that contain~$G_0$.
\end{Rem}

\begin{Cor}
\label{corconind}
Let $G$ be a finitely generated group and
$\mu$ a positive measure on~$G$ with finite support such that $G_\mu=G$.
A $\mu$-harmonic character $\chi$ of $G$
is never induced from a proper subgroup $G_0\subset G$.
\end{Cor}

\begin{proof}
Since $G_\mu=G$, the restriction $\mu_0$ of $\mu$ to $G_0$ satisfies $\mu_0<\mu$. Since $\chi$ is a character, one has $P_{\mu_{_0}}\chi =\al \chi$ with some constant $\al>0$.
Since $P_\mu\chi=\chi$, one has $\al<1$.
Therefore, one has $\lim_{n\ra\infty}P^n_{\mu_{_0}}\chi=0$, and
the $\mu$-harmonic function $\chi$ is not induced
from $G_0$.
\end{proof}

\subsection{Double induction}
\label{secdouind}

\bq
The following lemma
tells us that two successive inductions
of a positive harmonic function is equivalent to a direct induction.
\eq
Let $G$ be a finitely generated group.

\begin{Lem}
\label{lemdouind}
Let
$\mu_0<\mu'_0<\mu$ be positive measures on $G$ with finite support.
Let $h_0$ be a positive $\mu_0$-harmonic function on $G$.
The following are equivalent:
\begin{enumeratei}
\item\label{lemdouindi}
the function
$h:=\lim_{n\ra\infty}P^n_{\mu}h_0$ is finite.
\item\label{lemdouindii}
the functions
$h'_0:=\lim_{n\ra\infty}P^n_{\mu'_{_0}}h_0$ and
$h':=\lim_{n\ra\infty}P^n_{\mu}h'_0$ are finite.
\end{enumeratei}
\noindent
In this case, the two induced $\mu$ harmonic functions are equal $h=h'$.
\end{Lem}

\skpt
\begin{proof}[Proof of Lemma \ref{lemdouind}]
$\eqref{lemdouindi}\Rightarrow\eqref{lemdouindii}$
Since $h_0\leq h$, one has the inequalities
$P^n_{\mu'_0}h_0\leq P^n_{\mu'_0}h\leq P^n_{\mu}h= h$ and $h'_0\leq h$.
Therefore, one also has the inequalities $P^n_{\mu}h'_0\leq P^n_{\mu}h= h$ and $h'\leq h$.

$\eqref{lemdouindii} \Rightarrow\eqref{lemdouindi}$
Since $h_0\leq h'_0$, one has $P^n_{\mu}h_0\leq P^n_{\mu}h'_0$ and $h\leq h'$.
\end{proof}

\subsection{Induction of characters}
\label{secindcha}

\bq
We give now a few conditions that have to be satisfied
in order for the induction of a harmonic character to be a finite function.
\eq

Let $G$ be a finitely generated group and
$\mu$ be a positive measure on $G$ with finite support $S$
such that $G=G_\mu$.
We write $\mu =\mu_0+\mu_1$ as a sum of two positive measures
and set $S_0:=\supp\mu_0$ and $G_0:=G_{\mu_{_0}}$.
Let $\chi_0$ be a $\mu_0$-harmonic character of $G_0$
that we extend by $0$ as a function on $G$.
We denote by
\[
Z_G(G_0):=\{g\in G\mid gg_0=g_0g\text{ for all $g_0$ in $G_0$}\}
\]
the centralizer of $G_0$ in $G$, and by
\[
N_G(G_0,\chi_0):=\{g\in G\mid gg_0g^{\!-1}\!\!\in\! G_0
\text{ and } \chi_0(gg_0g^{\!-1})=\chi_0(g_0)\text{ for all $g_0$ in $G_0$}\}
\]
the normalizer of $(G_0,\chi_0)$ in $G$.

\begin{Lem}
\label{lemindcha}
If the induced $\mu$-harmonic function $h_{G_{_0},\chi_{_0}}$ is finite, then:
\begin{enumeratei}
\item\label{lemindchai}
The measure $\mu_0$ is the restriction of $\mu$ to $G_0$ and $S_0=S\cap G_0$.
\item\label{lemindchaii}
The subgroup $G_0$ has infinite index in $G$.
\item\label{lemindchaiii}
One has $G_{\mu_{_1}}^+\cap G_0=\emptyset$.
\item\label{lemindchaiv}
One has $G_{\mu_{_1}}^+\cap Z_G(G_0)=\emptyset$.
\item\label{lemindchav}
One has $G_{\mu_{_1}}^+\cap N_G(G_0,\chi_0)=\emptyset$.
\end{enumeratei}
\end{Lem}

\skpt
\begin{Rem}
\begin{itemize}
\item
In particular, the supports $S_0$ of $\mu_0$ and $S_1$ of $\mu_1$ are disjoint
and the semigroup $G_{\mu_{_1}}^+$ does not meet the center $Z$ of $G$.

\item
Note also that if one wants $h_{G_{_0},\chi_{_0}}$ to be $\mu$-extremal, the group
$G_0$ must be generated by $S_0$. Indeed if this is not the case,
the $\mu_0$-harmonic character $\chi_0$ is not $\mu_0$-extremal and,
by Lemma \ref{lemrecind},
the function $h_{G_{_0},\chi_{_0}}$ is not $\mu$-extremal.

\item
The above conditions are not the only necessary conditions,
as we will see in Section \ref{secexiind}.

\end{itemize}
\end{Rem}

\skpt
\begin{proof}[Proof of Lemma \ref{lemindcha}]
\eqref{lemindchai} This is equivalent to $\mu_1(G_0)=0$ which follows from \eqref{lemindchaiii}.

\eqref{lemindchaii} This follows from \eqref{lemindchaiii}. Indeed pick an element $s_1$
in the support of $\mu_1$, if the index were finite, there would exist a
positive power
$s_1^d$ belonging to $G_0$.

\eqref{lemindchaiii} This follows from \eqref{lemindchav} because $G_0\subset N_G(G_0,\chi_0)$.

\eqref{lemindchaiv} This follows from \eqref{lemindchav} because $Z_G(G_0)\subset N_G(G_0,\chi_0)$.

\eqref{lemindchav} This point is the main content of Lemma \ref{lemindcha}. We proceed by contraposition.
Let $S_1$ be the support of $\mu_1$ and
$w_1=s_1\ldots s_\ell\in S_1^\ell$, with $\ell\geq 1$ be a word
such that~$\dot{w}_{1}$ belongs to $N_G(G_0,\chi_0)$.

The proof relies on a cautious analysis of the words that occur in Equality \eqref{eqnpmun}.
We~recall the notation
$\mu_{1,w_{_1}}:=\mu_{1,s_{_1}}\!\cdots\mu_{1,s_{_\ell}}>0$.
We will denote $P_{w_{_1}}$ for the operator
of left translation by $\dot{w}_{1}:=s_1\cdots s_\ell\in G$; it is given by
$P_{w_{_1}}h(g)=h(\dot{w}_{1}g)$ for all function $h$ on $G$ and all $g$ in $G$.
One computes
\begin{align*}
P_\mu^{n+\ell}\chi_0(\dot{w}_{1}^{-1})
&\geq \sum_{1\leq i\leq n}
\mu_{1,w_{_1}} P^i_{\mu_{_0}}P_{w_{_1}}P^{n-i}_{\mu_{_0}}\chi_0(\dot{w}_{1}^{-1})\\
&= \sum_{1\leq i\leq n}
\mu_{1,w_{_1}} P^i_{\mu_{_0}}P_{w_{_1}}\chi_0(\dot{w}_{1}^{-1})
\hspace{2em}\text{because $\chi_0$ is $\mu_0$-harmonic}\\
&= \sum_{1\leq i\leq n} \mu_{1,w_{_1}}
\sum_{w_{_0}\in S_0^i}\mu_{0,w_{_0}}
\chi_0(\dot{w}_{1}\dot{w}_{0}\dot{w}_{1}^{-1})
\hspace{1em}\text{by definition of $P_{\mu_{_0}}$}\\
&= \sum_{1\leq i\leq n} \mu_{1,w_{_1}}
\sum_{w_{_0}\in S_0^i}\mu_{0,w_{_0}}
\chi_0(\dot{w}_{0})
\hspace{1em}\text{because $\dot{w}_{1}$ normalizes $\chi_0$}\\
&= \sum_{1\leq i\leq n} \mu_{1,w_{_1}}
\chi_0(0)= n\mu_{1,w_{_1}}
\hspace{1em}\text{because $\chi_0$ is $\mu_0$-harmonic.}
\end{align*}
This goes to infinity with $n$, and the induced function is not finite.
\end{proof}

\subsection{Negligible trajectories}
\label{secnegtra}
\bq
We now discuss a lemma on non-induced extre\-mal positive $\mu$-harmonic functions.
This lemma will be useful for the proof of the $Z$\nobreakdash-semiinvariance
of these functions on the Heisenberg group.
\eq

Let $G$ be a finitely generated group
and $\mu$ be a positive measure on $G$ with finite support $S$
generating $G$.

For every word $w=s_1\ldots s_n\in S^n$, we define
$k_w\geq 0 $ to be the smallest integer $k$ for which we can
write $w=w_0\ldots w_k$ as a concatenation of \emph{strongly non-generating} subwords $w_j$.
\emph{Strongly non-generating} means that
there exists an infinite index subgroup $G_j$ of $G$
containing all the letters $s_i$ occurring
in the subword $w_j$.
The following lemma tells us that the words with $k_w$ bounded
are negligible
in the sum \eqref{eqnpmun} for a non-induced $\mu$-harmonic function.

\begin{Lem}
\label{lemnegtra}
Let $h$ be a non-induced positive $\mu$-harmonic function on $G$
Then, for~all $k\geq 0$, and $g$ in $G$, the partial sums
\begin{equation}
\label{eqnlimsumhgg}
I_{n,k}(g):=\sum_{\substack{w\in S^n\\ k_w\leq k}} \mu_w\, h(\dot{w}g)
\end{equation}
converge to $0$ when $n\ra\infty$.
\end{Lem}

\begin{proof}[Proof of Lemma \ref{lemnegtra}]
Fix $g$ in $G$. For $w$ in $S^n$ we introduce
the maximal \emph{strongly non-generating} suffix $\si$ of $w$.
Suffix means that one can write $w=w'\si$.
We denote by $S_{0,w}$ the set of letters of $\si$ and
by $\ell_{0,w}$ the length of $\si$.
Since there are only finitely many subsets $S_0$ of $S$,
we can write $I_{n,k}(g)$
as a finite sum $I_{n,k}(g)=\sum I_{n,k,S_{_0}}(g)$,
where $I_{n,k,S_{_0}}(g)$ involves the words $w$ for which $S_{0,w}=S_0$.
Here this finite sum is indexed by the subsets $S_0$ of $S$ that generates an infinite index subgroup of $G$.
We~argue by induction on $k$.

\subsubsection*{First assume $k=0$}
For such $S_0\subset S$
one has
\[
I_{n,0,S_{_0}}(g)\leq
\sum_{w_{_0}\in S_{_0}^{n}}\mu_{w_{_0}}\, h(\dot{w}_0g)
=
P_{\mu_{_0}}^nh(g),
\]
where $\mu_0$ is the
restriction of $\mu$ to $S_0$.
By Definitions \ref{defnonind} and \ref{defnonindsuite}, since $h$ is non-induced
and since $S_0$ generates an infinite index subgroup of $G$, the sequence
$P_{\mu_{_0}}^nh(g)$ converges to $0$ when $n\ra \infty$,
and the claim \eqref{eqnlimsumhgg} is true for $k=0$.

\subsubsection*{Now assume $k\geq 1$}
Fix $\eps_0>0$.
Since $h$ is non-induced, as above, we can choose~$\ell_0$ such that,
for any subset $S_0$ of $S$ that generates an infinite index subgroup of $G$,
one has $P_{\mu_{_0}}^{\ell_0}h(g)\leq \eps_0$, where $\mu_0$ is the
restriction of $\mu$ to $S_0$.
We decompose the sum $I_{n,k,S_{_0}}(g)$ as a sum of two terms
\[
I_{n,k,S_{_0}}(g)=
I'_{n,k,S_{_0},\ell_{_0}}(g)+I''_{n,k,S_{_0},\ell_{_0}}(g),
\]
where
$I'_{n,k,S_{_0},\ell_{_0}}(g)$ involves the words $w$ for which $\ell_{0,w}\geq \ell_0$ and
$I''_{n,k,S_{_0},\ell_{_0}}(g)$
involves the words $w$ for which $\ell_{0,w}< \ell_0$.

Bounding $I'_n$. One computes, using the $\mu$-harmonicity of $h$,
\begin{align*}
I'_{n,k,S_{_0},\ell_{_0}}(g)&\leq
\sum_{w_{_0}\in S_{_0}^{\ell_0}}\mu_{w_{_0}}
\sum_{w_{_1}\in S^{n\!-\!\ell_0}}\,
\mu_{w_{_1}}\, h(\dot{w}_{1}\dot{w}_0g)\\
&\leq
\sum_{w_{_0}\in S_{_0}^{\ell_0}}\mu_{w_{_0}}\, h(\dot{w}_0g)\leq \eps_0.
\end{align*}

Bounding $I''_n$. One decomposes $I''_{n,k,S_{_0},\ell_{_0}}$
as a finite sum
\[
I''_{n,k,S_{_0},\ell_{_0}}(g)=\sum_{\si} \mu_{\si}I''_{n,k,\si}(g)
\]
over the finitely many words $\si$ of length $\ell<\ell_0$,
where
\begin{align*}
I''_{n,k,\si}(g)& \leq
\sum_{w'\in S^{n-\ell},\, k_{w'}\leq k-1}\, \mu_{w'}\, h({\dot{w}}'\dot{\si}g)\\
&\leq I_{n-\ell,k-1}(\dot{\si}g).
\end{align*}
Therefore by the induction hypothesis one has $\lim_{n\ra\infty}I''_{n,k,\si}(g)=0$.
Since $\eps_0$ can be chosen arbitrarily small, one deduces that
$\lim_{n\ra\infty}I_{n,k}(g)=0$.
\end{proof}

\subsection{When $G_\mu^+$ meets the center}
\label{secgmucen}

\bq
There is a simple case where the semiinvariance of $\mu$-harmonic
functions is easy to prove, namely when $G_\mu^+$ meets the center.
\eq

Let $G$ be a finitely generated group, $Z$ be the center of $G$ and
$\mu$ be a finite positive measure on $G$.

\begin{Lem}
\label{lemgmucen}
Assume that an element $z$ of $Z$
belongs to the semigroup $G_\mu^+$.
Then, for every extremal positive
$\mu$-harmonic function $h$ on $G$
there exists a constant $q>0$ such that $h_z=qh$.
\end{Lem}

We recall that $h_z$ is the function $g\mto h(gz)$.

\begin{proof}[Proof of Lemma \ref{lemgmucen}]
This is a slight generalization of the Choquet--Deny
theorem. Let $n\geq 1$ be an integer such that $z$ is in the support
of $\mu^{*n}$. The equality $h=P^n_\mu h$ is of the form
$h= \al h_z +h'$, where $\al>0$ and $h'$ is a positive function.
Since the function~$h_z$ is also $\mu$-harmonic, the extremality of $h$
implies that $h_z$ is proportional to $h$.
\end{proof}

\section{\texorpdfstring{$Z$}{Z}-Invariance of harmonic functions}
\label{secinvhar}

In all this section we keep the following notation:
\begin{equation}
\label{eqnnotsta}
\begin{array}{l}
\textit{$G$ is a finite index subgroup in $H_3(\m Z)$,
$Z$ is the center of $G$,}\\
\textit{$\mu$ is a positive measure with finite support $S$
such that $G_\mu=G$,}\\
\textit{$h$ is a positive $\mu$-harmonic function on $G$.}
\end{array}
\end{equation}
In this section we will mainly focus on non-induced $\mu$-harmonic functions
(see Definitions \ref{defnonind} and \ref{defnonindsuite})
and we will prove that they are $Z$-invariant.

We begin by a lemma that explain our choices in Definition \ref{defnonindsuite}.

\begin{Lem}
\label{lemnonind}
The positive $\mu$-harmonic function $h$ is non-induced if and only if
$\lim_{n\ra\infty}P_{\mu_0}^nh=0,$
for all restriction $\mu_0$ of $\mu$ to an abelian subset $S_0$ of $S$.
\end{Lem}

\begin{proof} By Definition \ref{defnonindsuite}, ``$h$ non-induced'' means that $h$
is non-induced from an infinite index subgroup $G_0$ of $G$.
Note that the subgroups $G_0$ of infinite index in $G$ are exactly
the abelian subgroups of $G$. Indeed any two non-commuting elements of $H_3(\m Z)$
generate a finite index subgroup of $H_3(\m Z)$.
\end{proof}

\begin{Rem}

A finite index subgroup $G$ of $H_3(\m Z)$
is not always isomorphic to $H_3(\m Z)$,
but it contains a finite index subgroup
that is isomorphic to $H_3(\m Z)$.
Extending our theorem \ref{thharmu1} to these groups $G$
would be straightforward but not so interesting.

The main reason we want to work with this slightly larger class of group $G$
in this section
is that, in the ``proof by induction'' of Proposition \ref{provmunol},
we need to apply
the ``induction hypothesis'' to a finite index subgroup of $G$.
\end{Rem}

\subsection{Semiinvariance of harmonic functions}
\label{secseminv}

\bq
In this section we prove that $h$
is semiinvariant by one central element.
The proofs below are self-contained.
They are inspired by the more intuitive
proofs for the south-west measure in \cite{BenoistHeisenberg1}
that rely on Young diagrams.
\eq

\begin{Prop}

\label{proseminv}
Keep notation \eqref{eqnnotsta} and
assume that $h$ is $\mu$-extremal and non-induced.
Then there exist $z\neq 0$ in $Z$ and $q>0$ such that $h_z=qh$.
\end{Prop}

\subsubsection*{Proof of Proposition \ref{proseminv}}
By Lemma \ref{lemgmucen}, we can assume $S\cap Z=\emptyset$.

For $n\geq 2$, we introduce a symmetric relation on $S^n$ given by
\begin{multline*}
\mc R_n:=\{ (w,w')\in S^n\times S^n \mid
w=w_0ss'w'_0\text{ and } w'=w_0s'sw'_0,\text{ where}\\
w_0\in S^i,\;w'_0\in S^{n-i-2},\; s\in S,\;s'\in S\text{ with } ss'\neq s's\}.
\end{multline*}
This means that $w$ and $w'$ are obtained from one another by switching two
consecutive non-commuting letters.
For a word $w\in S^n$ we let
\[
\text{$k_w$ $=$ the number of
pairs of consecutive non-commuting letters in $w$.}
\]
Since $G$ is the Heisenberg group $H_3(\m Z)$ and since $S\cap Z=\emptyset$,
this number
$k_w$ is the same as the one occurring in Lemma \ref{lemnegtra}.
Indeed, there exists a unique partition
$S=S_0\cup\ldots\cup S_\ell$ of $S$ such that two elements $s$, $s'$ of $S$
commute if and only if they belong to the same $S_i$.
To go on the proof of Proposition \ref{proseminv}, we will need the following two lemmas.

We set $p_0:=\max\limits_{s,s'\in S} |c_{s,s'}|$, where the integers $c_{s,s'}$
are defined in \eqref{eqncggzoc}.

\begin{Lem}
\label{lemkwkw}
For $(w,w')\in \mc R_n$, one has
\begin{enumeratei}
\item\label{lemkwkwi}
$\dot{w}=\dot{w}'z_0^p$ for some integer $p$ with $0<|p|\leq p_0$,
\item\label{lemkwkwii}
$\mu_{w'}=\mu_w$,
\item\label{lemkwkwiii}
$|k_{w'}-k_w|\leq 2$.
\end{enumeratei}
\end{Lem}

\skpt
\begin{proof}[Proof of Lemma \ref{lemkwkw}]
\eqref{lemkwkwi} This follows from the equality $ss'=s's\, z_0^{c_{s,s'}}$.

\eqref{lemkwkwii} The same letters occur in $w$ and $w'$.

\eqref{lemkwkwiii} The pairs of adjacent letters in $w$ and $w'$ are the same except for at most two of them.
\end{proof}

\begin{Lem}
\label{lemhghgzp}
For $g$ in $G$, one has $h(g)\leq \sum_{0<|p|\leq p_0}h(z_0^pg)$.
\end{Lem}

\begin{proof}[Proof of Lemma \ref{lemhghgzp}]
Replacing $h$ by its translate $h_g$, we can assume that $g=0$.
We want to prove that the following difference is non-positive:
\[
D:=
h(0)- \sum_{0<|p|\leq p_0}h(z_0^p) \leq 0.
\]
Using notations \eqref{eqnpmun}, we compute $D$ as
\[
D=
\sum_{w\in S^n}\mu_wh(\dot{w})-
\sum_{0<|p|\leq p_0}\sum_{w'\in S^n}\mu_{w'}h(\dot{w}'z_0^p).
\]
We fix $\eps_0>0$ and $k_0\geq 2+2\eps_0^{-1}$.
By Lemma \ref{lemnegtra}, one can find an integer $n\geq 1$
such that
the first sum limited at the trajectories $w$ for which $k_w<k_0$ is bounded by $\eps_0$. Using the fact that, for $w$ in $S^n$, the fiber
\[
\{(w,w')\mid w'\in S^n,\; (w,w')\in \mc R_n\}
\]
of the maps
$\mc R_n\mto S^n;\,(w,w')\mto w$ has cardinality $k_w$, one gets
\[
D\leq
\eps_0+
\sum_{\substack{(w,w')\in \mc R_{_n},\\ k_w\geq k_0}}
\Bigl(\frac{\mu_w}{k_w}\,h(\dot{w})-
\frac{\mu_{w'}}{k_{w'}}
\textstyle\sum_{0<|p|\leq p_0}h(\dot{w}'z_0^p)\Bigr).
\]
By Lemma \ref{lemkwkw}, the element $\dot{w}$ is equal
to at least one of those $\dot{w}'z_0^p$, therefore one gets
\[
D\leq
\eps_0+
\sum_{\substack{(w,w')\in \mc R_{_n},\\ k_w\geq k_0}}
\mu_w\frac{k_{w'}-k_w}{k_wk_{w'}}\,h(\dot{w}).
\]
By Lemma \ref{lemkwkw}, one has $|k_{w'}\!-\!k_w|\leq 2$,
and $2/k_{w'}\leq 2/(k_0\!-\!2)\leq \eps_0$, and
\[
D\leq
\eps_0+
\eps_0
\sum_{(w,w')\in \mc R_{_n}}
\frac{\mu_w}{k_w}\,h(\dot{w}).
\]
Using again that $k_w$ is the cardinality of the fiber and using the harmonicity of $h$,
one gets
\[
D\leq
\eps_0+
\eps_0
\sum_{w\in S^n}
\mu_w h(\dot{w})
=
\eps_0+
\eps_0\, h(0).
\]
Since $\eps_0$ can be chosen arbitrarily small, this gives $D\leq 0$ as expected.
\end{proof}

\begin{proof}[End of proof of Proposition \ref{proseminv}]
Lemma \ref{lemhghgzp} tells us that there exists a finite subset $F\subset Z\smallsetminus\{0\}$ and a positive $\mu$-harmonic function $h'$ such that
\[
\sum_{z\in F} h_z=h+h'.
\]
Since the cone $\mc H_\mu^+$ is well-capped and reticulated,
both the function $h'$ and the sum $\sum_{z\in F}h_z$ admit a unique desintegration in $\mu$-extremal functions
(see Section~\ref{secconhar}).
Hence, since all the positive $\mu$-harmonic functions $h$ and $h_z$ are $\mu$-extremal,
the function $h$ has to be proportional to one of these translates $h_z$.
\end{proof}

\begin{Rem}
We now want to deduce from the semi-invariance of $h$
proved in Proposition \ref{proseminv}, the $Z$-invariance of $h$.
This is not a general fact.
Indeed, the harmonic function $h$ in Case \eqref{thconclu3b} of Theorem \ref{thconclu}
can be $Z$-semiinvariant but is not $Z$-invariant.
Hence, we have to use once more the assumption that $h$ is not induced.
One technical difficulty comes from the fact that, when
$G^+_\mu\neq G$, the cone $\mc H_\mu^+$
often does not have a compact basis. This prevents us from using
the same arguments as in \cite{Margulis66}.
\end{Rem}

\subsection{$z$-invariance and $Z$-invariance}
\label{seczinzin}

\bq
We first notice that in order to prove the $Z$-invariance
of a positive $\mu$-harmonic function $h$ on the Heisenberg group $G$, it is enough
to check that it is invariant under one non trivial element of $Z$.
\eq

\begin{Lem}
\label{lemzinzin}
Keep notation \eqref{eqnnotsta} and
assume that there exists $z\neq 0$ in $Z$ such that
$h_z=h$. Then $h$ is $Z$-invariant.
In particular, if $h$ is $\mu$-extremal, it is proportional to a $\mu$-harmonic character of $G$.
\end{Lem}

Note that in this lemma the positive $\mu$-harmonic function $h$
is not assumed to be $\mu$-extremal.

\begin{proof}[Proof of Lemma \ref{lemzinzin}]
We write $z=z_0^p$. We can assume that $p$
is the smallest positive integer for which $h_z=h$.
We can also assume that $h$ is extremal in the
convex cone
\[
\mc H_{\mu,z}^+:=\{
\text{positive, $\mu$-harmonic and $z$-invariant functions on $G$}\}.
\]
Therefore the functions $h_{z_{_0}^i}$, for $i=1,\ldots, p$,
are non-proportional functions which are extremal
in this cone,
and the function $f:= h_{z_{_0}} +\cdots + h_{z_{_0}^{p}}$ is
$\mu$-harmonic and $Z$-invariant.

We claim that $f$ is extremal
among the $\mu$-harmonic functions on $G/Z$.
Indeed, assume that one can write $f=f'+f''$ with both $f'$ and $f''$
positive, $\mu$-harmonic and $Z$-invariant.
We argue as in the proof of Proposition \ref{proseminv}
with the well-capped and reticulated cone $\mc H_{\mu,z}^+$.
Both the function $f'$ and $f$
admit a unique desintegration in extremal functions in this cone
(see Section \ref{secconhar}).
Hence, since all the functions $h_z$ are extremal in this cone,
one must have $f'=\sum_{1\leq i\leq p} \la_i h_{z_{_0}^i}$
for some constants $\la_i\geq 0$. Since $f'$ is $z_0$-invariant,
all these constants are equal to some $\la\geq 0$ and one has $f'=\la f$.
This proves that $f$ is extremal
among the $\mu$-harmonic functions on $G/Z$.

Since $G/Z$ is abelian, by the Choquet--Deny theorem, this function $f$
is a $\mu$\nobreakdash-har\-monic character of $G$.
Therefore, by Lemma \ref{lemcarext}, this function $f$ is $\mu$-extremal
and one has $p=1$. This means that $h$ is $Z$-invariant.
\end{proof}

\subsection{$Z$-invariance when $V^+_\mu$ contains a line}
\label{secvmulin}

\bq
In this section, we finish the proof of our main Theorem
\ref{thharmu1}
when the cone $V_\mu^+$ is the plane or a half-plane,
see \eqref{eqnvmuphc}.
\eq

\begin{Prop}

\label{provmulin1}
Keep notation \eqref{eqnnotsta}, assume that
the cone $V^+_\mu$ contains a line and
that the $\mu$-harmonic function $h$ is not induced. Then
$h$ is $Z$-invariant.
\end{Prop}

\begin{proof}[Proof of Proposition \ref{provmulin1}]
We can assume that $h$ is $\mu$-extremal
and apply Proposition \ref{proseminv}.
Then our claim follows from the following slightly stronger
Proposition \ref{provmulin2}.
This stronger version will also be useful in Section \ref{secexiind}.
\end{proof}

\begin{Prop}

\label{provmulin2}
Keep notation \eqref{eqnnotsta} and assume that
the cone $V^+_\mu$ contains a line.
Assume also that there exists $z\neq 0$ in $Z$ and $q>0$ such that
$h_z=qh$. Then the function $h$ is $Z$-invariant.
\end{Prop}

\begin{proof}[Proof of Proposition \ref{provmulin2}]
According to Lemma \ref{lemzinzin}, it is enough to prove that $q=1$.
Replacing $h$ by a multiple of a suitable translate, we can assume that $h(0)= 1$.
Replacing $z$ by its inverse if necessary, we can also assume that $q\geq 1$.
Since the cone~$V_\mu^+$ contains a line,
there exists two words
$w_0$ in $S^{n_0}$ and $w'_0$ in $S^{n'_0}$
whose product is in the center:
\[
\dot{w}_0\dot{w}'_0= z^a
\quad\text{for some $a$ in $\m Z$.}
\]
Since the cone $V^+_\mu$ is not a line, there exists also a word $w_1$ in $S^{n_1}$ such that
\begin{equation}
\label{eqnbilcom0}
\dot{w}_0\dot{w}_1\dot{w}^{-1}_0\dot{w}_1^{-1}= z^b
\quad\text{for some $b\geq 1$.}
\end{equation}
Note that one might have to switch $w_0$ and $w'_0$
to ensure that $b\geq 1$.

Assume, for a contradiction, that $q\neq 1$, so that $q>1$. Choose an integer $\ell\geq 1$
such that $C:=\mu_{w_{_0}}\mu_{w'_{_0}}\,q^{a+b\ell}>1$. Notice the equality, for all $k\geq 1$,
\begin{equation}
\label{eqnbilcom1}
\dot{w}_0^k \,\dot{w}_1^\ell\, {\mbox{$\dot{w}$}'_0}^k\, \dot{w}_1^{-\ell} = z^{ak+b\ell k}.
\end{equation}
Note that both Equations \eqref{eqnbilcom0} and \eqref{eqnbilcom1} rely on
the bilinear formula \eqref{eqncggzoc}
for the commutators in the Heisenberg group.
Now we can compute with $n:=kn_0\!+\!\ell n_1\!+\! kn'_0$,
\begin{align*}
\label{eqnhwlcln}
h(\dot{w}_1^{-\ell})&=
P_\mu^{n}h(\dot{w}_1^{-\ell})\\
&\geq
\mu_{w_{_0}}^k \mu_{w_{_1}}^\ell \mu_{w'_{_0}}^k\,
h(\dot{w}_0^k, \dot{w}_1^\ell\, {\mbox{$\dot{w}$}'_0}^k\, \dot{w}_1^{-\ell})\\
&\geq
\mu_{w_{_0}}^k \mu_{w_{_1}}^\ell \mu_{w'_{_0}}^k\, q^{ak+b\ell k}
= \mu_{w_{_1}}^\ell C^k.
\end{align*}
Since $C>1$ and since this inequality is valid for all integer $k\geq 1$
one gets a contradiction. This proves that $q= 1$.
\end{proof}

\subsection{$Z$-invariance when $V^+_\mu$ contains no line}
\label{secvmunol}

\bq
In this section, we finish the proof of our main Theorem
\ref{thharmu1}
when the cone $V_\mu^+$ is properly convex, see \eqref{eqnvmuphc}.
\eq

\begin{Prop}

\label{provmunol}
Keep notation \eqref{eqnnotsta}, assume that
$V^+_\mu$ contains no line
and that the $\mu$-harmonic function $h$ is not induced.
Then $h$ is $Z$-invariant.
\end{Prop}

\subsubsection*{Beginning of proof of Proposition \ref{provmunol}}
The proof is by induction on the cardinality of the support $S$ of $\mu$,
simultaneously for all the finite index subgroups $G$ of $H_3(\m Z)$.
We will use the induction hypothesis inside the proof of Lemma \ref{lempmupmu}. We begin the proof by a few reduction steps.

\subsubsection*{First step}
We can assume that $h$ is $\mu$-extremal.
Indeed by Definitions \ref{defnonind} and~\ref{defnonindsuite},
almost all the $\mu$-extremal $\mu$-harmonic positive functions $f$
that occur in the desintegration $h=\int_{\mc E}f\rd \al (f)$ of $h$
are non-induced.
In this case,
by Proposition \ref{proseminv}, there exist $z=z_0^p\neq 0$ in $Z$ and $q>0$
such that $h_z=qh$.
According to Lemma \ref{lemzinzin}, it is enough to prove that $q=1$.

\begin{enumeratea}
\item\label{enum:a}
\emph{We can assume $z=z_0$.}
Because we can replace $h$ by
the function $f:= q_0^{-1}h_{z_{_0}} +\cdots + q_0^{-p}h_{z_{_0}^{p}}$,
where $q_0>0$ is chosen so that $q_0^p=q$. This function $f$ is
$\mu$-harmonic and $Z$-semiinvariant. It might not be $\mu$-extremal,
but this property will not be used in the argument below.

\item
\emph{We can assume $S\cap Z=\emptyset$.} Indeed, by \eqref{enum:a}, if $\mu_Z$ is the restriction of
$\mu$ to the center, one has $P_{\mu_{_Z}}h=\la h$ for a constant $0\leq \la <1$.
But then the function $h$ is harmonic for the measure $(1-\la)^{-1}(\mu-\mu_Z)$.
It might not be extremal for this measure,
but, as we just said, this is not important.

\item
\emph{We can assume $h(0)= 1$.}
Because we can replace $h$ by a multiple of a suitable translate.

\item
\emph{We can assume $q< 1$.} Because we can replace the generator $z_0$
by its inverse. We are now looking for a contradiction.
\end{enumeratea}

\subsubsection*{Second step}
We now can enter the key argument of the proof.
Since the cone $V_\mu^+$ is properly convex and since $S\cap Z=\emptyset$,
we can find a partition of the support of $\mu$ in two non-empty subsets
\begin{equation}
\label{eqns1s2}
S=S_1\cup S_2,
\end{equation}
such that
\begin{equation}
\label{eqncss}
c_{s_1,s_2}\geq 1\quad
\text{for all $s_1$ in $S_1$ and $s_2$ in $S_2$.}
\end{equation}
The partition \eqref{eqns1s2} is given by a suitable decomposition
$V_\mu^+=V_1^+\cup V_2^+$ of the properly convex cone $V_\mu^+$
in two cones $V_1^+$ and $V_2^+$ of disjoint interior so that
the inequalities~\eqref{eqncss} will follow from the bilinear formula \eqref{eqncggzoc}
for the commutators in the Heisenberg group.

We will use the decomposition $\mu=\mu_1+\mu_2$, where $\mu_1:=\mathbf{1}_{S_{_1}}\,\mu$ and
where \hbox{$\mu_2:=\mathbf{1}_{S_{_2}}\,\mu$}.
The proof again starts with the equality \eqref{eqnpmun}
which tells us that, for all $n\geq 1$,
\begin{equation}
\label{eqnhgshwg2}
1=h(0)=
\sum_{w\in S^n}\, \mu_w h(\dot{w}).
\end{equation}
We cut this sum into pieces
parametrized by pairs $(w_1,w_2)\!\in\!S_1^{n_1}\!\times\!S_2^{n_2}$,
with \hbox{$n_1\!+\!n_2\!=\!n$}. We define
\[
B_{w_1,w_2}=\{ w\in S^n \text{ containing $w_1$ and $w_2$ as subwords}\}.
\]
For instance when
$w_1=11$ and $w_2=23$,
one has
\[
B_{w_1,w_2}=\{1123,1213,1231,2113,2131,2311\}.
\]
This allows us to write the above sum \eqref{eqnhgshwg2} as
\begin{equation}
\label{eqnhgshwg3}
1=
\sum_{n_{_1}+n_{_2}=n}\;
\sum_{w_{_1}\in S_{_1}^{n_1}}
\sum_{w_{_2}\in S_{_2}^{n_2}}\;
\sum_{w\in B_{w_{_1}\!,w_{_2}}}\, \mu_w h(\dot{w}).
\end{equation}
For every $w$ in $B_{w_{_1},w_{_2}}$, we write, using iteratively \eqref{eqncggzoc}
and \eqref{eqncss},
\begin{equation}
\label{eqngwgwgw}
\dot{w}=\dot{w}_{2}\dot{w}_{1}z_0^{n_w}
\quad\text{for some integer $n_w\geq 1$}.
\end{equation}
Then Equality \eqref{eqnhgshwg3} becomes
\begin{equation}
\label{eqnhgshwg4}
1=
\sum_{n_{_1}+n_{_2}=n}\;
\sum_{w_{_1}\in S_{_1}^{n_1}}
\sum_{w_{_2}\in S_{_2}^{n_2}}\;
\mu_{w_{_1}} \mu_{w_{_2}} h(\dot{w}_{2}\dot{w}_{1})\,
\biggl(\sum_{w\in B_{w_{_1}\!,w_{_2}}}\,q^{n_w}\biggr).
\end{equation}
To pursue our analysis, we will need the following lemma which
bounds this last sum.

\begin{Lem}
\label{lemsumbww}
For all $w_1$ in $S_1^{n_1}$ and $w_2$ in $S_2^{n_2}$, one has
\begin{equation}
\label{eqnswbwbw}
\sum_{w\in B_{w_{_1}\!,w_{_2}}} q^{n_w}
\leq
\eta(q)^{-1}<\infty,
\end{equation}
where $\eta(q):=\prod\limits_{i\geq 1}(1-q^i)>0$.
\end{Lem}

Note that this upper bound does not depend on $(w_1,w_2)$.

\begin{proof}[Proof of Lemma \ref{lemsumbww}]
For each word $w=s_1\ldots s_n$ in $B_{w_{_1}\!,w_{_2}}$, we set
\[
m_w:=\bigl|\{(i,j)\mid 1\leq i< j\leq n
\quad\text{and}\quad s_i\in S_1,\; s_j\in S_2\}\bigr|.
\]
Condition \eqref{eqncss} implies that
\[
m_w \leq n_w
\quad\text{for all $w$ in $B_{w_{_1}\!,w_{_2}}$}.
\]
A word $w=s_1\ldots s_n$ in $B_{w_{_1}\!,w_{_2}}$ is determined by the
increasing sequence $1\leq i_1<i_2<\ldots <i_{n_{_2}}\leq n$ of places $i$
where $s_i$ belongs to $S_2$, and $m_w$ is given by
\[
m_w=(i_{n_{_2}}\!-\!n_2)+\cdots +(i_2\!-\! 2)+(i_1\!-\! 1).
\]
Therefore, for all $m\geq 1$,
the number
\[
p(n_1,n_2,m):=|\{w \in B_{w_{_1},w_{_2}}\mid
m_w=m\}|
\]
is equal to the number of partitions of $m$
by $n_2$ non-increasing integers $a_1,\ldots,a_{n_{_2}}$ bounded by $n_1$:
\[
p(n_1,n_2,m)=\bigl|\{n_1\geq a_1\geq \ldots \geq a_{n_{_2}}\geq 0
\quad\text{and}\quad m=a_1+\cdots+a_{n_{_2}}\}\bigl|.
\]
This quantity is bounded by the partition function
\[
p(m)=\bigl|\{a_1\geq \ldots \geq a_k\geq \ldots\geq 0
\quad\text{and}\quad m=a_1+\cdots+a_k+\cdots\}\bigl|.
\]
The generating function of the partition function is
\[
\label{eqnspmqmp}
\sum_{m\geq 0}p(m)q^m
=
\prod\limits_{i> 0}(1+q^i+q^{2i}+\cdots)
= \prod\limits_{i> 0}(1-q^i)^{-1}
= \eta(q)^{-1}.
\]
We now collect the sequence of inequalities
we have just proved
\begin{align*}
\sum_{w\in B_{w_{_1}\!,w_{_2}}} q^{n_w}
&\leq
\sum_{w\in B_{w_{_1}\!,w_{_2}}} q^{m_w}
=
\sum_{m\geq 0} p(n_1,n_2,m)q^{m}\\
&\leq
\sum_{m\geq 0} p(m)q^{m}
=
\eta(q)^{-1}
\end{align*}
and we obtain the bound \eqref{eqnswbwbw} we were looking for.
\end{proof}

\begin{proof}[End of proof of Proposition \ref{provmunol}]
We plug Inequality \eqref{eqnswbwbw} in Formula \eqref{eqnhgshwg4}
and we obtain, for all $n\geq 1$
\begin{equation}
\label{eqnpmupmu}
\sum_{n_{_1}+n_{_2}=n}
P^{n_1}_{\mu_{_1}}P^{n_2}_{\mu_{_2}}h(0)\geq \eta(q)>0.
\end{equation}
This contradicts the following Lemma \ref{lempmupmu}
\end{proof}

\begin{Lem}\label{lempmupmu}
With the same notation. In particular $\mu=\mu_1+\mu_2$
with $S_1$ and $S_2$ disjoint,
and $h$ is a non-induced $\mu$-harmonic function on $G$.
\begin{enumeratea}
\item\label{lempmupmua}
One has
$\lim_{n\ra\infty}P^{n}_{\mu_{_1}}h =0$ and $\lim_{n\ra\infty}P^{n}_{\mu_{_2}}h = 0$.
\item\label{lempmupmub}
One also has
\begin{equation}
\label{eqnpmupmo}
\lim_{n\ra\infty}\sum_{n_{_1}+n_{_2}=n}
P^{n_1}_{\mu_{_1}}P^{n_2}_{\mu_{_2}}h=0.
\end{equation}
\end{enumeratea}
\end{Lem}

\skpt
\begin{proof}[Proof of Lemma \ref{lempmupmu}]
\eqref{lempmupmua}
Let us prove it for $\mu_1$.

\begin{itemize}
\item
If $S_1$ is abelian, this follows from the assumption that $h$ is non-induced.

\item
If $S_1$ is not abelian,
we will use our induction hypothesis. Assume, for a contradiction, that
the $\mu_1$-harmonic function $h':=\lim_{n\ra\infty}P^{n}_{\mu_{_1}}h$
is non-zero. By Lemma~\ref{lemrecind},
this function $h'$ is $\mu_1$-extremal and satisfies
\begin{equation}
\label{eqnlimpnh}
\lim_{n\ra\infty}P^{n}_{\mu}h'=h.
\end{equation}
By Lemma \ref{lemdouind}, this $\mu_1$-harmonic function $h'$ is not induced and,
since $S_1$ is smaller than $S$,
the function $h'$ is a $\mu_1$-harmonic character
of the group $G_{\mu_{_1}}$.
Since this group $G_{\mu_{_1}}$ has finite index in the group $G_\mu$, Lemma \ref{lemindcha}\,\eqref{lemindchaii}
tells us that the function $\lim_{n\ra\infty}P^{n}_{\mu}h'$ is not finite.
This contradicts \eqref{eqnlimpnh}.
\end{itemize}

\eqref{lempmupmub}
The argument is the same as for Lemma \ref{lemnegtra}, but is simpler.
We fix $g$ in $G$ and $\eps_0>0$.
According to point $a)$, there exists $N_1\geq 1$ such that
$P^{N_1}_{\mu_{_1}}h(g)\leq \eps_0$.
Let~$I_n$ be the left-hand side of \eqref{eqnpmupmo}.
We decompose $I_n(g)$
as the sum of two terms
\[
I_n(g)=I'_n(g)+I''_n(g),
\]
where $I'_n(g)$ involves the terms with $n_1\geq N_1$
and $I''_n(g)$ involves the terms with $n_1 < N_1$

Bounding $I'_n$. One computes, using the $\mu$-harmonicity of $h$,
\begin{align*}
I'_{n}(g)&=
\sum_{n'_{_1}+n_{_2}=n-N_{_1}}
P^{N_1}_{\mu_{_1}}P^{n'_1}_{\mu_{_1}}P^{n_{_2}}_{\mu_{_2}}h(g)\\
&\leq
P^{N_1}_{\mu_{_1}}P^{n-N_1}_{\mu}h(g)
= P^{N_1}_{\mu_{_1}}h(g)\leq \eps_0.
\end{align*}

Bounding $I''_n$. One decomposes $I''_n(g)$ as a finite sum
\[
I''_{n}(g)=\sum_{n_{_1}<N_{_1}}\sum_{w_{_1}\in S_{_1}^{n_1}} \mu_{w_{_1}}P^{n-n_1}_{\mu_{_2}}(\dot{w}_{1}g)
\]
over the finitely many words $w_1$ of length $n_1<N_1$.
By point $a)$, all terms of the sum go to $0$ so that one has
$\lim_{n\ra\infty}I''_{n}(g)=0$.

Since $\eps_0$ can be chosen arbitrarily small, one gets
$\lim_{n\ra\infty}I_{n}(g)=0$.
\end{proof}

This ends the proof of Proposition \ref{provmunol}. We can now complete the proof of our main theorem \ref{thharmu1}.

\begin{proof}[Proof of Theorem \ref{thharmu1}.]
Let $h$ be an extremal positive $\mu$-harmonic function on $G$.
By Propositions \ref{provmulin1} and \ref{provmunol}, either $h$ is $Z$-invariant or
$h$ is induced.

Assume first that $h$ is $Z$-invariant, then $h$ is an extremal positive harmonic function
on the abelian group $G/Z$ and, by the Choquet--Deny theorem (see Section \ref{secpreres}),
$h$ is proportional to a character of $G$.

Assume now that $h$ is induced. Since $h$ is extremal,
as we have seen in Lemma~\ref{lemrecind} and Definitions \ref{defnonind} and \ref{defnonindsuite},
there exist an infinite index subgroup
$G_0$ of $G$ and an extremal $\mu_0$-harmonic function on $G_0$,
where $\mu_0$ is the restriction of $\mu$ to $G_0$,
such that the function $h$ is a translate of the function $h_{G_{_0},h_{_0}}$ induced from $h_0$.
Since $G$ is the Heisenberg group, this group $G_0$ is abelian
and, by the Choquet--Deny theorem, the extremal $\mu_0$-harmonic function $h_0$ is proportional to a character of $G_0$.
\end{proof}

\section{Existence of induced harmonic functions}
\label{secexiind}

In this section, except for Section \ref{secrankfour}, we will keep the following notations:
\begin{equation}
\label{eqnnotbul}
\begin{array}{l}
\textit{$G=H_3(\m Z)$ is the Heisenberg group,
$Z$ is the center of $G$,}\\
\textit{$\mu$ is a positive measure with finite support $S$
such that $G_\mu=G$,}\\
\textit{$S_0\subset S$ is a maximal abelian subset, $\mu_0:=\mathbf{1}_{S_{_0}}\mu$,
$G_0:=G_{\mu_{_0}}$,}\\
\textit{$\chi_0$ is a $\mu_0$-harmonic character of $G_0$ and $\mu_1:=\mu-\mu_0$.}
\end{array}
\end{equation}
By Theorem \ref{thharmu1}, we know that an extremal positive $\mu$-harmonic functions
on~$G$ which is not proportional to a character is proportional to
a translate of an induced $\mu$\nobreakdash-harmonic function of the form $h_{G_0,\chi_0}$.
Note that the maximality of $S_0$ is guaranteed by Lemma \ref{lemindcha}\,\eqref{lemindchaiv}.

We will give in this section a
necessary and sufficient condition
for the induced $\mu$-harmonic function $h_{G_0,\chi_0}$
to be finite.

In Lemma \ref{lemindcha}\,\eqref{lemindchaiv}, we have already found that the following condition
is necessary: $G_{\mu_{_1}}^+\cap Z_G(G_0)=\emptyset$.
Since $V_{\mu_{_0}}$ is a line, one can check that this condition is equivalent~to:
\begin{equation}
\label{eqnvmuvmu}
S_1\cap Z=\emptyset
\quad\text{and}\quad
V^+_{\mu_{_1}}\cap V_{\mu_{_0}}=\{0\}.
\end{equation}
We will assume that it is satisfied. We distinguish two cases according to the rank of the abelian group $G_0$.

\subsection{Induction of characters when $\rank G_0=1$}
\label{secindhar1}
\bq
In this section we give the necessary and sufficient condition for
the induced function $h_{G_{_0},\chi_{_0}}$ to be finite when the rank of $G_0$ is $1$.
\eq

Note that, since $S_0$ is maximal abelian in $S$,
one has the equivalence:
\[
\rank G_0=1\Longleftrightarrow G_0\cap Z=\{0\}.
\]

\begin{Prop}
\label{proindra1}
Keep notation \eqref{eqnnotbul}. Assume \eqref{eqnvmuvmu} and $\rank G_0=1$.
Then the induced harmonic function $h:=h_{G_{_0},\chi_{_0}}$ is finite if and only if
the probability measure $\widetilde{\mu}_0:=\chi_0\mu_0$ on $G_0$ is not centered.
\end{Prop}

\skpt
\begin{Rem}
\begin{itemize}
\item
The measure $\widetilde{\mu}_0=\chi_0\mu_0$ is a probability measure because
$\chi_0$ is a $\mu_0$-harmonic character.
\item
The condition {\it $\widetilde{\mu}_0$ centered} means, as usual, that
$\sum_{s\in S_{_0}}\widetilde{\mu}_{0,s}\, \ol s =0$ in $V$,
where $\ol{s}$ is the image of $s$ in $V$.
\item
This condition {\it $\widetilde{\mu}_0$ non-centered} is always satisfied
when $V_\mu^+$ contains no line.
\end{itemize}
\end{Rem}

\begin{proof}[Proof of Proposition \ref{proindra1}]
Using \eqref{eqnvmuvmu} and $\rank G_0=1$, we can assume that
\[
S_0\subset \{ (x,0,0)\mid x\in \m Z\}
\quad\text{and}\quad
S_1\subset \{ (x,y,z)\in G\mid y\geq 1\}.
\]
Let $\tau:G_0\mto \m Z$ be the morphism given by
$\tau(g_0)=x$ for $g_0=(x,0,0)$.

\subsubsection*{First case: when $\widetilde{\mu}_0$ is centered}
We fix $s_1$ in $S_1$ and we compute, as in Lemma \ref{lemindcha}, for $n\geq 1$,
\begin{equation}
\label{eqnhspnch}
\begin{aligned}
h(s_1^{-1})&\geq
P_\mu^{n+1}\chi_0(s_1^{-1})
\geq
\mu_{s_{_1}}\sum_{k\leq n}
P^k_{\mu_{_0}}P_{s_1}P^{n-k}_{\mu_{_0}}\chi_0(s_1^{-1})\\
&=
\mu_{s_{_1}}\sum_{k\leq n}
P^k_{\mu_{_0}}P_{s_1}\chi_0(s_1^{-1})=
\mu_{s_{_1}}\sum_{k\leq n}
\sum_{w\in S_{_0}^k}\mu_{0,w}\,\chi_0(s_1\dot{w}s_1^{-1}).
\end{aligned}
\end{equation}
The words $w$ that contribute to this sum are those with
$s_1\dot{w}s_1^{-1}\dot{w}^{-1}\in G_0$, \ie
$\dot{w} =0$ or, equivalently, $\tau(\dot{w})=0.$
Hence letting $n$ go to $\infty$, one gets
\[
h(s_1^{-1})
\geq
\mu_{s_{_1}}\sum_{k\geq 0}
\sum_{w\in S_{_0}^k}\widetilde{\mu}_{0,w}\mathbf{1}_{\{\tau(\dot{w})=0\}}.
\]
If we write $w=s_1\ldots s_n$ and $x_i:=\tau(s_i)$,
and if we think of these letters $s_i$ as independent random variables with same law $\widetilde{\mu}_0$,
this inequality can be rewritten as
\[
h(s_1^{-1})
\geq
\mu_{s_{_1}}\sum_{k\geq 0}
\m P(x_1+\cdots +x_k=0).
\]
But since the random variables $x_i\in \m Z$ are centered,
the expected number of passage at $0$ of the walk $x_1+\cdots +x_k$ is infinite, and the function $h$ is not finite.

\subsubsection*{Second case: when $\widetilde{\mu}_0$ is not centered}
The computation is similar but more involved since
we want to prove finiteness of $h(g)$ at every point $g$ in $G$.

We want a uniform upper bound for
\[
P^n_\mu\chi_0(g)
=
\sum_{w\in S^n}\mu_w\,\chi_0(\dot{w}g).
\]
The only words $w$ that contribute to this sum are those for which
$\dot{w}g$ is in $G_0$.
By~assumption \eqref{eqnvmuvmu},
if we extract from $w$ the maximal subword $\si=s_1\ldots s_\ell$
whose letters are in $S_1$, the length $\ell$ of $\si$ is uniformly bounded
by an integer $\ell_0$.
Therefore we can split the above sum into a finite sum
\[
P^n_\mu\chi_0(g)
=
\sum_{\ell\leq \ell_0}\sum_{\si\in S_{_1}^\ell}
\mu_\si\,Q_{\si,n-\ell}\,\chi_0(g),
\]
where
\begin{equation}
\label{eqnqsinch}
\begin{aligned}
Q_{\si,n}\,\chi_0(g)
&=
\sum_{k_0+\cdots +k_\ell=n}
P^{k_\ell}_{\mu_{_0}}P_{s_{_\ell}}\cdots
P^{k_1}_{\mu_{_0}}P_{s_{_1}}P^{k_0}_{\mu_{_0}}\,\chi_0(g)\\
&=
\sum_{k_1+\cdots +k_\ell\leq n}
P^{k_\ell}_{\mu_{_0}}P_{s_{_\ell}}\cdots
P^{k_1}_{\mu_{_0}}P_{s_{_1}}\,\chi_0(g).
\end{aligned}
\end{equation}
We want to bound the limit
\begin{equation}
\label{eqnqsiinf}
\begin{aligned}
Q_\infty(g)
&:=
\lim_{n\ra\infty}
Q_{\si,n}\,\chi_0(g)\\
&\hphantom{:}=
\sum_{k_{_1}\geq 0}\cdots\sum_{k_{_\ell}\geq 0}
\sum_{w_{_1}\in S^{k_1}_{_0}}\cdots\sum_{w_{_\ell}\in S^{k_\ell}_{_0}}
\mu_{w_{_1}}\cdots\mu_{w_{_\ell}}\,\chi_0(s_1\dot{w}_{1}\cdots s_\ell\, \dot{w}_{\ell}\, g).
\end{aligned}
\end{equation}
For $i\leq \ell$, let $\si_i:=s_1\cdots s_i\in G$ and $b_i\geq 1$ be the integer given by
\[
\si_ig_0\si_i^{-1}g_0^{-1}={z_0}^{-b_i\tau(g_0)}
\quad\text{for all $g_0$ in $G_0$,}
\]
so that one has
\begin{equation}
\label{eqnswswg}
s_1\dot{w}_{1}\cdots s_\ell\, \dot{w}_{\ell}\, g
=
\dot{w}_{1}\cdots \dot{w}_{\ell}\,\si_\ell\, g\, {z_0}^{-b_1\tau(\dot{w}_1)-\cdots -b_\ell\tau(\dot{w}_\ell)}.
\end{equation}
Writing $\si_\ell\, g=g_0\,z_0^{c}$ with $g_0$ in $G_0$ and $c$ in $\m Z$
one gets
\[
%\label{eqnqinf2}
Q_\infty(g)
=\chi_0(g_0)\sum_{k_{_1}\geq 0}\ldots\sum_{k_{_\ell}\geq 0}
\sum_{w_{_1}\in S^{k_1}_{_0}}\ldots\sum_{w_{_\ell}\in S^{k_\ell}_{_0}}
\widetilde{\mu}_{0,w_{_1}}\cdots\widetilde{\mu}_{0,w_{_\ell}}\,
\mathbf{1}_{\{b_1\tau(\dot{w}_1)+\cdots +b_\ell\tau(\dot{w}_\ell)=c\}}.
\]
If we think of all the letters occurring in one of the words $w_1,\ldots, w_\ell$
as independent random variables with same law $\widetilde{\mu}_0$,
this equality can be written~as
\[
\label{eqnqinf3}
Q_\infty(g)
=\chi_0(g_0)
\sum_{k_{_1}\geq 0}\ldots\sum_{k_{_\ell}\geq 0}
\m P(b_1S_{1,k_{_1}}+\cdots +b_\ell S_{\ell,k_{_\ell}}=c),
\]
where $S_{i,k_{_i}}:=\tau(\dot{w}_i)$.
Then the finiteness of $Q_\infty(g)$ follows from the following
Lemma~\ref{lemranwal}.
\end{proof}

\begin{Lem}
\label{lemranwal}
Let $(X_{i,k})_{i\leq \ell,\, k\geq 1}$, be independent real variables
with same law. Assume this law has finite support and is not centered.
Let $(b_i)_{i\leq\ell}$ be positive numbers and $c$ be a real number.
Set $S_{i,k}:=X_{i,1}+\cdots+X_{i,k}$.
Then one has
\begin{equation}
\label{eqnranwal}
\sum_{k_{_1}\geq 0}\cdots\sum_{k_{_\ell}\geq 0}
\m P(b_1S_{1,k_{_1}}+\cdots +b_\ell S_{\ell,k_{_\ell}}=c)
<\infty.
\end{equation}
\end{Lem}

\begin{proof}[Proof of Lemma \ref{lemranwal}]
We adapt the classical proof of the large deviation inequality.
We set $X=X_{1,1}$. Assume for instance that $\m E (X)>0$.
One can choose $\eps>0$ so that all the expectations
\[
\al_i:=\m E(e^{-\eps b_i X})
\]
are smaller than $1$. Then one computes
\begin{align*}
\m P(b_1S_{1,k_{_1}}+\cdots +b_\ell S_{\ell,k_{_\ell}}=c)
&\leq
\m E(e^{\eps\,(c -b_1S_{1,k_{_1}}-\cdots -b_\ell S_{\ell,k_{_\ell}})})\\
&=
e^{\eps c}\,\m E(e^{-\eps b_1X})^{k_1}\cdots \m E(e^{-\eps b_\ell X})^{k_\ell}=
e^{\eps c}\,\al_1^{k_1}\cdots \al_\ell^{k_\ell}
\end{align*}
and therefore, summing all these inequalities, we find the following upper bound
for the left-hand side $L$ of \eqref{eqnranwal}
\[
L\leq e^{\eps c}\,(1-\al_1)^{-1}\cdots (1-\al_\ell)^{-1} <\infty.
\]
This ends the proof of the lemma and of Proposition \ref{proindra1}.
\end{proof}

\subsection{Induction of characters when $\rank G_0=2$}
\label{secindhar2}
\bq
In this section we give the necessary and sufficient condition for
the induced function $h_{G_{_0},\chi_{_0}}$ to be finite when
the rank of $G_0$ is $2$,
or equivalently when $G_0\cap Z\neq\{0\}$.
\eq

We split the statement into two cases depending
on the shape of the convex cone~$V^+_\mu$.

\begin{Prop}
\label{proindra2}
Keep notation \eqref{eqnnotbul}. Assume \eqref{eqnvmuvmu} and $\rank G_0=2$.
Assume moreover that the cone $V^+_\mu$ contains a line. Then
the induced harmonic function $h:=h_{G_{_0},\chi_{_0}}$ is not finite.
\end{Prop}

\begin{proof}[Proof of Proposition \ref{proindra2}]
This follows from Proposition \ref{provmulin2}.
Indeed, let $z$ be a non-zero element of $G_0\cap Z$ and $q:=\chi_0(z)$.
Assume, for a contradiction, that the function~$h$ is finite.
By Lemmas \ref{lemcarext} and \ref{lemconind},
this function $h$ is $\mu$-extremal.
By construction this function $h$ is semiinvariant: one has
$h_z=qh$. Hence by Proposition \ref{provmulin2}, one has $q=1$
and by Lemma \ref{lemzinzin} the $\mu$-harmonic function $h$ is $Z$-invariant.
Therefore, by~the Choquet--Deny theorem, this function $h$
is a $\mu$-harmonic character of $G$.
But by Corollary \ref{corconind}, a $\mu$-harmonic character is never induced.
Contradiction.
\end{proof}

\begin{Prop}

\label{proindra3}
Keep notation \eqref{eqnnotbul}. Assume \eqref{eqnvmuvmu} and $\rank G_0=2$.
Assume moreover that the cone $V^+_\mu$ contains no line.
Then the induced harmonic function $h:=h_{G_{_0},\chi_{_0}}$ is finite if and only if
\begin{equation}
\label{eqnsssmsm}
\text{there exist $s_0$ in $S_0$ and $s_1$ in $S_1$ such that $\chi_0(s_0s_1s_0^{-1}s_1^{-1})>1$.}
\end{equation}
\end{Prop}

\begin{Rem} Even though we will not use this remark, it is interesting to notice
that, since Assumption \eqref{eqnvmuvmu} is satisfied and since
the cone $V^+_\mu$ is properly convex,
this condition \eqref{eqnsssmsm} is equivalent to\vspace*{-3pt}
\begin{equation}
\label{eqnsssmsm2}
\textit{for all $s_0$ in $S_0\smallsetminus Z$ and $s_1$ in $S_1$ one has $\chi_0(s_0s_1s_0^{-1}s_1^{-1})>1$.}
\end{equation}
\end{Rem}

\begin{proof}[Proof of Proposition \ref{proindra3}]
The calculation is the same as for Proposition \ref{proindra1},
but the interpretation is different.
Using \eqref{eqnvmuvmu} and the proper convexity of the cone $V_\mu^+$,
we~can assume that\vspace*{-3pt}
\begin{equation}
\label{eqnSOS1}
S_0\subset \{ (x,0,z)\in G\mid x\geq 0\}
\quad\text{and}\quad
S_1\subset \{ (x,y,z)\in G\mid y\geq 1\}.
\end{equation}
Let $\tau:G_0\mto \m Z$ be the morphism given by
$\tau(g_0)=x$ for $g_0=(x,0,z)$.

\subsubsection*{Proof of $\Rightarrow$}
By \eqref{eqnSOS1}, we know that the half-line $V^+_{\mu_{_0}}$ is extremal
in the properly convex cone $V^+_\mu$.
Assume by contraposition, that
for all $s_0$ in $S_0$ and $s_1$ in $S_1$
one has $\chi_0(s_0s_1s_0^{-1}s_1^{-1})\leq 1$.
In particular, one has\vspace*{-3pt}
\begin{equation}
\label{eqnchi0chi0}
\chi_0(s_1\dot{w}s_1^{-1})
\geq \chi_0(\dot{w})
\quad\text{for all $s_1\in S_1$ and $w\in S_0^k$.}
\end{equation}

We fix $s_1$ in $S_1$ and, using \eqref{eqnchi0chi0}, we compute, for $n\geq 1$, as in \eqref{eqnhspnch},\vspace*{-3pt}
\begin{align*}
h(s_1^{-1})
&\geq
\mu_{s_{_1}}\sum_{k\leq n}
\sum_{w\in S_{_0}^k}\mu_{0,w}\,\chi_0(s_1\dot{w}s_1^{-1})\\[-3pt]
&\geq
\mu_{s_{_1}}\sum_{k\leq n}
\sum_{w\in S_{_0}^k}\mu_{0,w}\,\chi_0(\dot{w})\geq
\mu_{s_{_1}}\sum_{k\leq n}\chi_0(0)
\geq n\, \mu_{s_{_1}}.
\end{align*}
Letting $n$ go to $\infty$, one gets
$h(s_1^{-1})=\infty$.

\subsubsection*{Proof of $\Leftarrow$}
As for Proposition \ref{proindra1},
one can find an integer $\ell_0$ and one can split $P^n_\mu\chi_0(g)$
as a sum parametrized by words $\si=s_1\ldots s_\ell$ with letters in $S_1$
and $\ell\leq \ell_0$:\vspace*{-3pt}
\begin{align*}
P^n_\mu\chi_0(g)
&=
\sum_{\ell\leq \ell_0}\sum_{\si\in S_{_1}^\ell}
\mu_\si\,Q_{\si,n-\ell}\,\chi_0(g),
\quad\text{where, as in \eqref{eqnqsinch},}\\[-3pt]
Q_{\si,n}\,\chi_0(g)
&=
\sum_{k_1+\cdots +k_\ell\leq n}
P^{k_\ell}_{\mu_{_0}}P_{s_{_\ell}}\cdots
P^{k_1}_{\mu_{_0}}P_{s_{_1}}\,\chi_0(g).
\end{align*}
As in \eqref{eqnqsiinf}, we want to bound the limit\vspace*{-3pt}
\[
Q_\infty(g)
:=
\lim_{n\ra\infty}
Q_{\si,n}\,\chi_0(g).
\]
The only words $w$ that contribute to this sum are those for which
$\dot{w}g$ is in $G_0$.
For~$i\leq \ell$, let $\si_i:=s_1\cdots s_i\in G$ and $b_i\geq 1$ be the integer given by\vspace*{-3pt}
\[
\si_i\,g_0\,\si_i^{-1}g_0^{-1}=z_0^{-b_i\tau(g_0)}
\quad\text{for all $g_0$ in $G_0$,}
\]
so that one has\vspace*{-3pt}
\begin{equation}
\label{eqnswswg2}
s_1\dot{w}_{1}\cdots s_\ell\, \dot{w}_{\ell}\, g
=
\si_1\dot{w}_{1}\si_1^{-1}\cdots\, \si_{\ell}\,\dot{w}_{\ell}\,\si_\ell^{-1}\,\si_\ell\, g.
\end{equation}
Hence, one gets\vspace*{-3pt}
\[
Q_\infty(g)
=\mu_\si\,\chi_0(\si_\ell\, g)\,F_1\cdots F_\ell,
\]
where, for all $i\leq \ell$,\vspace*{-3pt}
\[
F_i:=
\sum_{k\geq 0}
\sum_{w\in S^{k}_{_0}}
\mu_{0,w}\, \chi_0(\si_i\, \dot{w}\,\si^{-1}_i).
\]
We want to prove that the sums $F_i$ are finite.
We will denote by $q_0>0$ the real number such that
for all $i$ in $\m Z$ such that $z_0^i$ is in $G_0$, one has
$\chi_0(z_0^i)=q_0^i.$
By assumption, one has $q_0>1$. One computes then\vspace*{-3pt}
\[
F_i
=
\sum_{k\geq 0}
\sum_{w\in S^{k}_{_0}}
\widetilde{\mu}_{0,w}\, q_0^{-b_i\tau(\dot{w})},
\]
where, as before, $\widetilde{\mu}_0$ is the probability
measure $\chi_0\mu_0$.
Let $p_w$ be the number of letters of $w$ that belong to $S_0\smallsetminus Z$
and $\al:=\widetilde{\mu}_0(S_0\smallsetminus Z)<1$.
One goes on:\vspace*{-3pt}
\begin{align*}
F_i
\leq
\sum_{k\geq 0}
\sum_{w\in S^{k}_{_0}}
\widetilde{\mu}_{0,w}\, q_0^{-p_w}
&=\sum_{k\geq 0}
\sum_{j\leq k}
\binom{j}{k}\al^j(1-\al )^{k-j}q_0^{-j}\\[-3pt]
&=
\sum_{k\geq 0}
(1-\al+\al q_0^{-1})^k
= \al^{-1}(1-q_0^{-1})^{-1}< \infty.
\end{align*}
This proves the finiteness of $F_i$, of $Q_\infty(g)$ and of the function $h_{G_{_0},\chi_{_0}}$.
\end{proof}

\subsection{Existence of harmonic characters}
\label{secharcha}

\bq
We explain in this section when $\mu_0$\nobreakdash-har\-monic characters
on abelian groups do exist.
\eq

Let $G_0=\m Z^d$ and $\mu_0$ be a
positive measure with finite support $S_0$ generating $G_0$ as a group.
For a character $\chi_0$ of $G_0$ we set
\[
\m E(\chi_0):=\sum_{s\in S_{_0}}\mu_{0,s}\chi_0(s).
\]
The map $\chi_0\ra\m E(\chi_0)$ is the Laplace transform of $\mu_0$.
We denote by
\begin{equation}
\label{eqnlamuoo}
\la(\mu_0):=\inf_{\chi_{_0}} \m E(\chi_0)
\end{equation}
the minimum of this Laplace transform.
Here is an example where it is easy to compute $\la(\mu_0)$.

\begin{Rem}
\label{remlammuo}
If $S_0$ is included in a properly convex cone of $\m R^d$, one has \hbox{$\la(\mu_0)=\mu_0(0)$}.
More generally, if $S_0$ is included in a half-space bounded by a hyperplane $H_0$, one has $\la(\mu_0)=\la(\mu_0|_{H_{_0}})$.
\end{Rem}

\skpt
\begin{Lem}
\label{lemharcha}
\begin{enumeratea}
\item\label{lemharchaa}
There exists a $\mu_0$-harmonic character if and only if
$\la(\mu_0)\leq 1$.

\item\label{lemharchab}
We can choose it so that $\widetilde{\mu}_0:=\chi_0\mu_0$ is not centered if and only if
$\la(\mu_0)< 1$.
\end{enumeratea}
\end{Lem}

\begin{proof} Lemma \ref{lemharcha} follows from the following three remarks:
\begin{itemize}
\item
A character $\chi_0$ is $\mu$-harmonic if and only if $\m E(\chi_0)=1$.\item
The group of characters is isomorphic to $\m R^d$, hence it is connected.
\item
Since $S_0$ contains non-zero elements one has $\sup_{\chi_{_0}}\m E(\chi_0)=\infty$.\qedhere
\end{itemize}
\end{proof}

\skpt
\begin{Cor}
\label{corharcha}
\begin{enumeratea}
\item
If $\mu_0(S_0)\leq 1$, $\mu_0$-harmonic characters exist.
\item
If $\mu_0(S_0)< 1$, we can choose it so that
$\widetilde{\mu}_0:=\chi_0\mu_0$ is not centered.
\item
If $\mu_0(S_0)>1$ and $\mu_0$ is centered, $\mu_0$-harmonic characters do not exist.
\end{enumeratea}
\end{Cor}

\begin{proof}
This follows from Lemma \ref{lemharcha} and the inequality $\la(\mu_0)\leq \mu_0(S_0)$.
\end{proof}

\subsection{Conclusion}
\label{secconclu}
\bq
We sum up in the following theorem
the main results we have obtained in this paper.
\eq

Let $G=H_3(\m Z)$ be the Heisenberg group,
$Z$ be the center of $G$, $\mu$ be a positive measure on $G$ with finite support $S$
such that $G_\mu=G$. We use the notation of Section~\ref{secheigro}.

\begin{Thm}
\label{thconclu}
The extremal positive $\mu$-harmonic functions on $G$ are
proportional either to a character
of $G$ or to a translate of a function $h$ induced from a character on an abelian subgroup.
Here is the list when $\mu(G)\!=\!1$.

\begin{enumerate}
\item\label{thconclu1}
\emph{When $V^+_\mu$ is the plane}. There is no induced $\mu$-harmonic function.

\item\label{thconclu2}
\emph{When $V^+_\mu$ is a half-plane}. Let $V_0$ be the boundary line of $V^+_\mu$
and $G_0\subset G$ be the subgroup
generated by the
elements of $S$ above $V_0$ and $\mu_0:=\mu|_{G_{_0}}$.
Then $h$ is equal to a function $h_{G_{_0},\chi_{_0}}$ induced from a $\mu_0$-harmonic character $\chi_0$ of $G_0$.
\begin{enumerate}
\item\label{thconclu2a}
If $G_0\cap Z = \{ 0\}$
there are exactly two such $h_{G_{_0},\chi_{_0}}$.

\item\label{thconclu2b}
If $G_0\cap Z\neq \{ 0\}$ there is no such $h_{G_{_0},\chi_{_0}}$.
\end{enumerate}

\item\label{thconclu3}
\emph{When $V^+_\mu$ is properly convex}. Let $V^+_i$, $i=0$, $1$, be the two
extremal rays of~$V^+_\mu$, let $G_i\subset G$,
be the two subgroups generated by the
elements of $S$ above $V^+_i$ and $\mu_i:=\mu|_{G_{_i}}$.
Then $h$ is equal to a function $h_{G_{_i},\chi_{_i}}$ induced from a $\mu_i$-harmonic character~$\chi_i$ of $G_i$, $i=0$ or $1$.
\begin{enumerate}
\item\label{thconclu3a}
If $G_i\cap Z = \{ 0\}$
there is exactly one such $h_{G_{_i},\chi_{_i}}$.

\item\label{thconclu3b}
If $G_i\cap Z\neq \{ 0\}$
there are uncountably many such $h_{G_{_i},\chi_{_i}}$.
\end{enumerate}
\end{enumerate}
\end{Thm}

\begin{Rem}
Theorem \ref{thconclu} is illustrated in the schematic Figures \ref{figcase1}, \ref{figcase2} and~\ref{figcase3} of the introduction. We have drawn a rough approximation
of the shape of the semigroup $G^+_\mu\subset G$
and its subsemigroups $G^+_{\mu_{_0}}$ and $G^+_{\mu_{_1}}$,
in order to illustrate the different cases that occur in Theorem \ref{thconclu}.
In these pictures the center $Z$ is the vertical axis.
\end{Rem}

\begin{proof}[Proof of Theorem \ref{thconclu}]
The first claim follows from
Proposition \ref{provmulin1} and Proposition \ref{provmunol}. Moreover, Case \eqref{thconclu1}
and
the first claims of Cases \eqref{thconclu2} and \eqref{thconclu3} follow from Lemma \ref{lemindcha}\,\eqref{lemindchaiv}.
\begin{itemize}
\item
\emph{Case} \eqref{thconclu2a} Since $\rank G_0=1$,
by Proposition \ref{proindra1}, $\chi_0$ must be a $\mu_0$-harmonic character
of $G_0$ with $\chi_0\mu_0$ non centered.
Since $\mu_0(G_0)<1$ and since $\mu_0$ is not supported by a half-line,
there are exactly two such $\chi_0$.
\item
\emph{Case} \eqref{thconclu2b} Since $\rank G_0=2$, this follows from Proposition \ref{proindra2}.
\item
\emph{Case} \eqref{thconclu3a} Since $\rank G_i=1$,
by Proposition \ref{proindra1},
$\chi_i$ must be a $\mu_i$-harmonic character of $G_i$
with $\chi_i\mu_i$ non centered.
Since $\mu_i(G_i)<1$ and since $\mu_i$ is supported by a half-line,
there is exactly one such $\chi_i$.
\item
\emph{Case} \eqref{thconclu3b} Since $\rank G_i=2$,
by Proposition \ref{proindra3}, $\chi_i$
must be a $\mu_i$-harmonic character of $G_i$
satisfying \eqref{eqnsssmsm}.
Since $\mu(G_i)<1$, and since $\mu_i$ is supported
by a half-space delimited by $Z$,
there are uncountably many such $\chi_i$.\qedhere
\end{itemize}
\end{proof}

\enlargethispage{-\baselineskip}%
\pagebreak[2]
\skpt
\begin{Rem}
\label{remconclu}
\begin{itemize}
\item
When $\mu$ is not assumed to be a probability measure,
the formulation of Theorem \ref{thconclu} has to be modified.
Indeed, if $\mu (\{ 0\})\geq 1$,
positive $\mu$-harmonic functions cannot exist.
More precisely, each of the three cases \eqref{thconclu2a}, \eqref{thconclu3a} and \eqref{thconclu3b}
has to be split into two subcases:
\begin{enumerate}\itemindent0pt
\item[(2a$'$)] If $G_0\cap Z = \{ 0\}$ and $\la(\mu_0)< 1$,
there are exactly two such $h_{G_{_0},\chi_{_0}}$. 
\item[(2a$''$)] If $G_0\cap Z = \{ 0\}$ and $\la(\mu_0)\geq 1$,
there are no such $h_{G_{_0},\chi_{_0}}$.
\item[(3a$'$)] If $G_i\cap Z = \{ 0\}$ and $\mu(\{ 0\})< 1$
there is exactly one such $h_{G_{_i},\chi_{_i}}$.
\item[(3a$''$)] If $G_i\cap Z = \{ 0\}$ and $\mu(\{ 0\})\geq 1$
there are no such $h_{G_{_i},\chi_{_i}}$.
\item[(3b$'$)] If $G_i\cap Z\neq \{ 0\}$
and
$\la(\mu_{_Z})_i\!<\! 1$,\!
there are uncountably many $h_{G_{_i},\chi_{_i}}$.
\item[(3b$''$)] If $G_i\cap Z \neq \{ 0\}$ and
$\la(\mu_{_Z})_i\geq 1$,
there are no such $h_{G_{_i},\chi_{_i}}$.
\end{enumerate}
\item
Here is the definition,
motivated by Condition \eqref{eqnsssmsm}, of the constants $\la(\mu_{_Z})_i$.
For~instance, for $i=0$, we choose $s_j\in \supp(\mu_j)\smallsetminus Z$,
$j=0$ or $1$,
and set\vspace*{-3pt}\enlargethispage{.5\baselineskip}%
\[\textstyle
\la(\mu_{_Z})_{_0}:=\inf\bigl\{\sum_{s\in Z}\mu_s\chi_{_Z}(s)\mid \chi_{_Z}
\text{ character of $Z$, $\chi_{_Z}(s_0s_1s_0^{-1}s_1^{-1})>1$}\bigr\}.
\]
\item
Note that Cases (2a$''$), (3a$''$) and (3b$''$) do not occur
when $\mu$ is a probability measure.
\end{itemize}
\end{Rem}

We now can deduce from Theorem \ref{thconclu} the corollaries in the
introduction:

\begin{proof}[Proof of Corollary \ref{corharmu1}]
The support of $h_{G_{_0},\chi_{_0}}$
is the semigroup generated by $G_0$ and $S^{-1}$.
In the cases where a $\mu$-harmonic function
$h_{G_{_0},\chi_{_0}}$ induced from a character is finite,
by Lemma \ref{lemindcha}\,\eqref{lemindchaiv}, one has $V^+_{\mu_{_1}}\cap V_{\mu_{_0}}=\{0\}$,
and this semigroup is never equal to $G$.
\end{proof}

\begin{proof}[Proof of Corollary \ref{corharmu2}]
Both conditions \eqref{corharmu2i}, \eqref{corharmu2ii}
are true in Cases \eqref{thconclu1} and \eqref{thconclu2b}.
Both conditions
are not true in Cases \eqref{thconclu2a}, \eqref{thconclu3a} and \eqref{thconclu3b}.
\end{proof}

One also has the following variation of Corollary \ref{corharmu2}.

\begin{Cor}
\label{corharmu3}
Same notation and $\mu(G)=1$.
The following are equivalent:
\begin{enumeratei}
\item\label{corharmu3i}
Extremal positive $\mu$-harmonic functions are $Z^{p}$-semiinvariant
for a $p\geq 1$.
\item\label{corharmu3ii}
$G_\mu^+\cap Z\not\subset\{0\}$.
\end{enumeratei}
\end{Cor}

Condition \eqref{corharmu3i} means that there exist $q>0$ and a non-zero element $z$ in $Z$
such that $h_z=qh$.

\begin{proof}[Proof of Corollary \ref{corharmu3}]
Both conditions \eqref{corharmu3i}, \eqref{corharmu3ii}
are true in Cases \eqref{thconclu1}, \eqref{thconclu2b}~and~\eqref{thconclu3b}.
Both conditions
are not true in Cases \eqref{thconclu2a} and \eqref{thconclu3a}.
\end{proof}

\subsection{A nilpotent group of rank $4$}
\label{secrankfour}

\bq
In this section, we explain why Theorem \ref{thharmu1}
can not be extended to all nilpotent groups.
\eq

In this section $G=N_4(\m Z)$ will be the nilpotent group equal to the set $\m Z^4$
of quadruples seen as matrices
$(t,x,y,z):=\mbox{\scriptsize $\left(\!
\begin{array}{cccc} 1 &t&t^2/2&z \\
0 &1&t&y\\0&0&1&x\\0&0&0&1
\end{array}\!\right)$}.$ The product is
\begin{equation*}
\label{eqnprofour}
(t,x,y,z)\, (t',x',y',z')= (t+t',x+x',y+y'+tx',z+z'+ty'+\tfrac12 t^2x').
\end{equation*}
The center $Z$ of $G$ is generated by $z_0:=(0,0,0,1)$.
Let $\mu$ be the measure
\begin{equation*}
\label{eqnmu0}
\mu:=\frac12(\de_{a}+\de_{b}),
\quad\text{where $a:=(1,0,0,0)$ and $b=(0,1,0,0)$.}
\end{equation*}
A $\mu$-harmonic function $h$ on $G$ is a function that satisfies, for all $g$ in $G$,
\begin{align*}
\label{eqnhghaghbg}
2\,h(g)&= h(ag)+ h(bg)
\quad\text{or, equivalently,}\\
2\,h(t,x,y,z)&= h(t\!+\!1,x,y+x,z\!+\!y\!+\! x/2)+ h(t,x\!+\!1,y,z).
\end{align*}

We now construct
extremal positive $\mu$-harmonic functions on $G$. Fix a sequence $\si$ of rapidly increasing integers $1\leq \si_0\leq \si_1\leq \si_2\leq\cdots$.
We introduce the left-infinite word $w_\si$ in $a$ and $b$ of the form
\[
w_\si = \cdots a^{(\si_{\la})}ba^{(\si_{\la-1})}b\cdots ba^{(\si_1)}ba^{(\si_0)},
\]
where the notation $a^{(m)}$ means that the letter $a$ is repeated $m$-times.
For each $k\geq 0$ we denote by $w_{\si,k}$ the suffix of length $k$ of $w_\si$,
\ie the word given by the $k$ last letters of $w_\si$.
We introduce the functions on $G$
\[
\psi_\si:=\sum_{k\geq 0}2^k\mathbf{1}_{\dot{w}_{\si,k}}
\quad\text{and}\quad
h_\si:=\sup_{n\geq 1}P_\mu^n\psi_\si.
\]
As before, the dot means that we replace the word by its image in $G$.

\begin{Lem}
\label{lemnilpfour}
Let $G=N_4(\m Z)$ and $\mu$ and $\si$ be as above. Assume that
\begin{equation}
\label{eqnsilaincrease}
\si_{\la+1}\geq 2\, \si_{\la}^2
\quad\text{for all $\la\geq 0$}.
\end{equation}
\begin{enumeratea}
\item\label{lemnilpfoura}
The function $\psi_\si$ is subharmonic and the sequence $P_\mu^n\psi_\si$ is increasing.
\item\label{lemnilpfourb}
The limit $h_\si=\lim_{n\to \infty}P_\mu^n\psi_\si$ is finite.
\item\label{lemnilpfourc}
The function $h_\si$ is an extremal positive $\mu$-harmonic function on $G$.
\item\label{lemnilpfourd}
The function $h_\si$ is not $Z$-invariant.
\item\label{lemnilpfoure}
The function $h_\si$ is not induced.
\end{enumeratea}
\end{Lem}

Note that Condition \eqref{eqnsilaincrease} is not optimized.

\skpt
\begin{proof}
\eqref{lemnilpfoura} One has $\psi_\si\leq P_\mu\psi_\si$ and therefore $P_\mu^n\psi_\si\leq P_\mu^{n+1}\psi_\si$.

\eqref{lemnilpfourb} This is the key point.
Fix $g=(t,x,y,z)\in G$.
Fix also two words $u$ and $v$ in $a$ and $b$
such that $g=\dot{u}\dot{v}^{-1}$. Such words always exist.
By definition of $h_\si$, one has
\begin{equation}
\label{eqnhsig}
h_{\si}(g)=2^{t+x} \lim_{\la\to\infty}
\text{(number of words $w$ such that $\dot{w}\dot{u}=\dot{w}_{\si,n_{_\la}}\dot{v}$),}
\end{equation}
where $w_{\si,n_{_\la}}$ is the suffix of $w_\si$
of length $n_\la:= \la+\sum_{ i\leq \la} \si_i$, \ie
\begin{equation}
\label{eqnwordwsi}
w_{\si,n_{_\la}} = a^{(\si_{\la})}ba^{(\si_{\la-1})}b\cdots ba^{(\si_1)}ba^{(\si_0)}.
\end{equation}
We also write
\begin{equation}
\label{eqnwordw}
w:= a^{(k_{\ell})}ba^{(k_{\ell-1})}b\cdots ba^{(k_1)}ba^{(k_0)},
\end{equation}
with all $k_i\geq 0$. We denote by $\ell(w)$ the length of a word $w$.

We want to prove, using Condition \eqref{eqnsilaincrease}, that the quantity \eqref{eqnhsig} is finite.
We~will see that there exists $\la_0=\la_0(g)$ such that
the number of words $w$ such that
\begin{equation}
\label{eqndotwgwsi}
\dot{w}\dot{u}=\dot{w}_{\si,n_{_\la}}\dot{v}
\end{equation}
does not depend on $\la$ for $\la\geq \la_0(g)$.
More precisely, we will see below that,
for $\la\geq \la_0(g)$,
Equality \eqref{eqndotwgwsi} implies
that $k_\ell=\si_{\la}$, so that we could remove the prefix $a^{(\si_\la)}b$
in both words and replace $\la$ by $\la\! -\! 1$.

Equality \eqref{eqndotwgwsi} gives four equations.
We could write them down but we will not need to.
The first two equations tell us that the same number of $a$'s and the same number of $b$'s
occur in the words $wu$ and $w_{\si,n_{_\la}}v$. In particular, those words have same length.
The third equation tells us that the sum of the positions of the $b$'s in these words
are the same.
Once these three equations are satisfied, one can write
$\dot{w}\dot{u}=\dot{w}_{\si,n_{_\la}}\dot{v} z_0^{N_w}$
with $N_w\in \m Z$.
The fourth equation
tells us that this integer $N_w$ is zero.

We claim that, for $\la\geq \la_0(g)$, if $k_\ell\neq \si_\la$ then $N_w\neq 0$.
The reason is that we can go from the word $wu$ to the word
$w_{\si,n_{_\la}}v$ by a (minimal) succession of ``moves'' that changes only the
central component.
These ``moves'' are
\[
w_1abw_2baw_3 \longleftrightarrow w_1baw_2abw_3.
\]
The images in $G$ are modified by a factor $z_0^{\ell_a(w_2)+1}$, where $\ell_a(w_2)$
is the number of~$a$'s occurring in the word $w_2$:
\[
\dot{w}_1ba\dot{w}_2ab\dot{w}_3 = \dot{w}_1ab\dot{w}_2ba\dot{w}_3z_0^{\ell_a(w_2)+1}.
\]

Therefore, by \eqref{eqnsilaincrease}, the largest contributions to $N_w$ come from the ``moves''
that involve the first $b$
on the left of the word $w_{\si,n_{_\la}}$. Hence,
when $k_\ell\neq \si_\la$ one has
\[
|N_{w}|\geq
\si_{\la-1}- (\ell(v)+
\sum_{i\leq \la-2}\si_i)^2
\]
which is non zero for $\la$ large enough by Condition \eqref{eqnsilaincrease}.

\eqref{lemnilpfourc} By construction $h_\si$ is $\mu$-harmonic.
We want to prove that $h_\si$ is extremal.
Assume that $h_\si=h'+h''$ with $h'$ and $h''$ positive $\mu$-harmonic.
It follows also from the previous computations
that
\[
2^{-k}h_\si(\dot{w}_{\si,k})=1
\quad\text{for all $k\geq 0$.}
\]
Introduce the two limits
\[
\al':=\lim_{k\ra\infty}2^{-k}h'(\dot{w}_{\si,k})
\quad\text{and}\quad
\al'':=\lim_{k\ra\infty}2^{-k}h''(\dot{w}_{\si,k}).
\]
These limits exist since by harmonicity of $h'$ and $h''$ these sequences are non-increasing.
Moreover, one has $\al'+\al''=1$ and
\[
h'\geq \al' \psi_\si
\quad\text{and}\quad
h''\geq \al'' \psi_\si.
\]
Using again the harmonicity of $h'$ and $h''$, one deduces
\[
h'\geq \al' h_\si
\quad\text{and}\quad
h''\geq \al'' h_\si.
\]
Since $\al'+\al''=1$, these inequalities must be equalities. This proves that $h$ is extremal.

\eqref{lemnilpfourd} The above computation also tells us that
$\supp(h_\si)\cap Z=\{0\}$.
This prevents~$h_\si$ to be $Z$-invariant.

\eqref{lemnilpfoure} If the function $h_\si$ were induced, it would be the translate by an element $g\in G$
of a function induced
from a character of the cyclic group $G_a$ generated by $a$ or
of the cyclic group $G_b$ generated by $b$.
Since $h_\si(w_k)\neq 0$, for all $k\geq 1$,
all the sets $G_\mu^+ w_k$ would meet $G_a g^{-1}$ or $G_b g^{-1}$, which is impossible since
both $\lim_{k\ra\infty}\ell_b(w_k)=+\infty$ and $\lim_{k\ra\infty}\ell_a(w_k)=+\infty$.
\end{proof}

\backmatter
